\documentclass{article}

\usepackage[margin=1in]{geometry}
\usepackage{graphicx}
\usepackage{booktabs}
\usepackage{longtable}
\usepackage{array}
\usepackage{tabularx}
\usepackage{listings}
\usepackage{xcolor}
\usepackage{amsmath}
\usepackage{hyperref}
\usepackage{float}
\usepackage[section]{placeins}
\usepackage{needspace}
\usepackage{subcaption}


\renewcommand{\arraystretch}{1.2}

\definecolor{codegray}{gray}{0.95}
\lstdefinestyle{cstyle}{
    language=C,
    backgroundcolor=\color{codegray},
    basicstyle=\ttfamily\scriptsize,
    keywordstyle=\color{blue!60!black}\bfseries,
    commentstyle=\color{green!35!black}\itshape,
    stringstyle=\color{orange!70!black},
    numbers=left,
    numberstyle=\tiny\color{gray},
    numbersep=6pt,
    breaklines=true,
    breakatwhitespace=false,
    showstringspaces=false,
    frame=single,
    tabsize=4,
    columns=fullflexible,
    captionpos=b,
}
\lstdefinestyle{makestyle}{
    language=make,
    backgroundcolor=\color{codegray},
    basicstyle=\ttfamily\scriptsize,
    breaklines=true,
    showstringspaces=false,
    frame=single,
    tabsize=4,
    columns=fullflexible,
    captionpos=b,
}
\lstset{style=cstyle}

\title{Predicting Cache Evolution: The Growing Capacity and Latency Cost of CPU Caching}
\author{Nick Ngo, Kevin Chen}
\date{August 2026}

\begin{document}

\maketitle

\section{Introduction}

The goal of this project is to reverse engineer the CPU cache hierarchy using controlled software microbenchmarks and timing measurements. The project compares cache behavior across multiple Intel, AMD, and Arm microarchitectures and CPU generations.

\section{Experimental Platforms and Architecture Research}

Table~\ref{tab:machines} summarizes the ECE/Ajorpaz lab machines used during the timing-only phase.

\begin{table}[htbp]
\centering
\small
\caption{ECE/Ajorpaz Lab Machines and Architecture Research}
\label{tab:machines}

\begin{tabularx}{\textwidth}{|l|X|l|l|l|l|l|}
\hline
\textbf{Machine} &
\textbf{ISA} &
\textbf{Vendor} &
\textbf{CPU Model} &
\textbf{Microarchitecture} &
\textbf{Year} &
\textbf{Node} \\
\hline

Sunbird &
x86-64 &
Intel &
2x Xeon E5-2680 v3 &
Haswell &
2014 &
22 nm\\
\hline

Thunderbird &
AArch64 &
Ampere &
Ampere Q80-30 &
ARM Neoverse N1 &
2020 &
7 nm\\
\hline

Skylark &
x86-64 &
AMD &
2x AMD EPYC 7532 &
Zen 2 &
2019 &
7 nm \\
\hline

Artemisia &
x86-64 &
Intel &
2x Xeon Gold 5420+ &
Sapphire Rapids &
2023 &
10 nm, Intel 7\\
\hline

Charnwood &
x86-64 &
Intel &
Core i7-6700 &
Skylake &
2015 &
14 nm\\
\hline

Crux &
x86-64 &
Intel &
Core i7-9700 &
Coffee Lake &
2019 &
14 nm\\
\hline

Ookay &
x86-64 &
Intel &
Core i7-7700 &
Kaby Lake &
2017 &
14 nm\\
\hline

Upgrade &
x86-64 &
Intel &
Core i7-8700 &
Coffee Lake &
2017 &
14 nm\\
\hline

\end{tabularx}

\end{table}

\section{Experimental Methodology}

\subsection{Part 1 -- Cache Benchmarks}

Every timing-only reverse-engineering target in this project -- cache-level/capacity
discovery, line size, associativity, hit latency, miss/next-level latency, and
inclusion/exclusion behavior (Section~\ref{sec:phase1}) -- is built from one shared
measurement primitive implemented once in \texttt{main\_code/common/} and compiled into a
single portable command-line tool, \texttt{cache\_bench}
(\texttt{main\_code/common/main.c}). Every experiment differs only in \emph{which}
parameter it sweeps and how it lays out the pointer-chase graph; the load, timing, and
batching machinery underneath is identical across all six targets and is described once
here so later subsections can point back to a listing number instead of repeating it, per
the reused-common-code convention required for this report.

\subsubsection{Dependent Pointer-Chase Primitive}

Every measurement in this codebase is built around a single dependent load chain: a
singly-linked cycle of fixed-size nodes, each holding only a \texttt{next} pointer, laid
out over a working-set buffer of controlled size and node spacing. Because the address of
load $i+1$ is the \emph{value} produced by load $i$, the CPU cannot issue, prefetch, or
reorder load $i+1$ before load $i$ completes -- this is what isolates true load-to-use
cache/TLB latency instead of memory bandwidth or issue-width effects. Listing~\ref{lst:pchase}
shows the whole primitive; it is architecture-independent and is reused, unmodified, by
every experiment in Section~\ref{sec:phase1}.

\begin{lstlisting}[caption={Dependent pointer-chase primitive (portable, all architectures). File: \texttt{main\_code/common/pointer\_chase.h} / \texttt{pointer\_chase.c}, commit \texttt{e0f2538}.},label={lst:pchase}]
#ifndef POINTER_CHASE_H
#define POINTER_CHASE_H

#include <stdint.h>

struct node {
    struct node *next;
};

void chase(struct node *p, uint64_t steps);

#endif

/* ---- pointer_chase.c ---- */
#include "pointer_chase.h"

static volatile struct node *sink_node;

void chase(struct node *p, uint64_t steps) {
    for (uint64_t i = 0; i < steps; i++) {
        p = p->next; /* next address depends on previous load */
    }
    sink_node = p; /* keep the final result observable */
}
\end{lstlisting}

The result of the final load is stashed in a \texttt{volatile} sink so the compiler cannot
prove the whole loop is dead and eliminate it -- necessary discipline even at the
\texttt{-O0} baseline used throughout this project (Section~\ref{sec:build}).

\subsubsection{Timing: Serialized Cycle Counters (x86-64 and AArch64)}

Every experiment measures elapsed time in raw CPU cycles ("ticks") via a portable
\texttt{timer\_start()}/\texttt{timer\_stop()} pair that dispatches, at compile time, to an
architecture-specific implementation (Listing~\ref{lst:timer}). This is the \emph{only}
architecture-specific code in the entire codebase -- every experiment-specific file
(Section~\ref{sec:phase1}) calls only the portable \texttt{timer.h} interface and requires
zero changes to run on Sunbird/Skylark/Crux/Ookay/Upgrade/Charnwood/ Artemisia (x86-64) or
Thunderbird (AArch64).

\begin{lstlisting}[caption={Portable timer dispatch. File: \texttt{main\_code/common/timer.h}, commit \texttt{e0f2538}.},label={lst:timer}]
#ifndef TIMER_H
#define TIMER_H

#include <stdint.h>

#if defined(__x86_64__) || defined(__i386__)

#include "../x86_64/timer_x86.h"

static inline uint64_t timer_start(void){
    return x86_tsc_start();
}

static inline uint64_t timer_stop(void){
    return x86_tsc_stop();
}

#elif defined(__aarch64__)

#include "../aarch64/timer_arm.h"

static inline uint64_t timer_start(void){
    return arm_counter_start();
}

static inline uint64_t timer_stop(void){
    return arm_counter_stop();
}

#else

#error "Unsupported architecture"

#endif

#endif
\end{lstlisting}

On x86-64 (Listing~\ref{lst:timerx86}), \texttt{timer\_start()} issues \texttt{LFENCE}
before \texttt{RDTSC} to drain in-flight instructions so the cycle read is not reordered
ahead of the code it is meant to time, and \texttt{timer\_stop()} uses the partially
self-serializing \texttt{RDTSCP} (which also drains prior instructions before reading the
counter) followed by a further \texttt{LFENCE} to prevent the \emph{next} region's
instructions from being reordered backward into the timed window.

\begin{lstlisting}[caption={x86-64 serialized TSC timer. File: \texttt{main\_code/x86\_64/timer\_x86.h}, commit \texttt{e0f2538}. Target architecture: x86-64.},label={lst:timerx86}]
#include <stdint.h>
#include <x86intrin.h>

static inline uint64_t x86_tsc_start(void) {
    _mm_lfence();
    uint64_t t = __rdtsc();
    _mm_lfence();
    return t;
}

static inline uint64_t x86_tsc_stop(void) {
    unsigned aux;
    uint64_t t = __rdtscp(&aux);
    _mm_lfence();
    return t;
}

static inline uint64_t time_one_load(volatile const uint64_t *p) {
    uint64_t t0 = x86_tsc_start();
    uint64_t value = *p;
    uint64_t t1 = x86_tsc_stop();
    __asm__ __volatile__("" : "+r"(value) :: "memory");
    return t1 - t0;
}
\end{lstlisting}

On AArch64 (Listing~\ref{lst:timerarm}), the free-running virtual counter
\texttt{CNTVCT\_EL0} is read between \texttt{DSB SY}/\texttt{ISB} barrier pairs, which
serialize the instruction stream in the same role \texttt{LFENCE} plays on x86. A real
resolution limitation surfaced on Thunderbird (Ampere Q80-30, Neoverse N1):
\texttt{CNTFRQ\_EL0} reports only 25\,MHz, so ticks/access for the whole 1\,KiB--1\,GiB
capacity range span roughly 0.08--2.4 -- below the x86-calibrated automatic boundary/knee
detectors' default absolute-tick thresholds (Sections~\ref{sec:capacity} and
\ref{sec:assoc}), which had to be rescaled per-machine on this ISA rather than reused
as-is.

\begin{lstlisting}[caption={AArch64 serialized virtual-counter timer. File: \texttt{main\_code/aarch64/timer\_arm.h}, commit \texttt{e0f2538}. Target architecture: AArch64.},label={lst:timerarm}]
#include <stdint.h>
static inline uint64_t arm_counter_start(void) {
    uint64_t t;
    __asm__ __volatile__(
        "dsb sy\n\t"
        "isb\n\t"
        "mrs %0, cntvct_el0\n\t"
        "isb\n\t"
        : "=r"(t) :: "memory");
    return t;
}

static inline uint64_t arm_counter_stop(void) {
    uint64_t t;
    __asm__ __volatile__(
        "dsb sy\n\t"
        "isb\n\t"
        "mrs %0, cntvct_el0\n\t"
        "isb\n\t"
        : "=r"(t) :: "memory");
return t;
}

static inline uint64_t arm_counter_freq(void) {
    uint64_t f;
    __asm__ __volatile__("mrs %0, cntfrq_el0" : "=r"(f));
    return f;
}
\end{lstlisting}

\subsubsection{Batched Measurement and Statistical Discipline}
\label{sec:batching}

A single \texttt{timer\_start()}/\texttt{timer\_stop()} pair carries tens of ticks of fixed
overhead (quantified directly in Section~\ref{sec:misslatency}), so every experiment except
miss/inclusion latency (which are inherently single-shot, Sections~\ref{sec:misslatency}
and~\ref{sec:inclusion}) times a \emph{batch} of $B$ back-to-back dependent accesses as one
region and reports the batch average, amortizing that fixed cost across $B{=}1000$
accesses by default. Batching this way -- rather than reporting one grand mean per
working-set size -- preserves the per-batch distribution, which downstream repeat-run and
box-plot analysis (Sections~\ref{sec:capacity}--\ref{sec:inclusion}) depends on to tell a
genuine cache-level transition apart from run-to-run system noise. The
\texttt{>= 1{,}000{,}000} timed-access-per-point requirement is enforced structurally, not
by convention: \texttt{num\_batches} is computed by ceiling division of
\texttt{samples}/\texttt{batch\_size}, and untimed warm-up passes never write into the
latency array that is later printed, so a warm-up access can never be miscounted as a timed
sample. Listing~\ref{lst:benchmark} also shows the required independent-load control used
by the hit-latency experiment (Section~\ref{sec:hitlatency}): each load's address is drawn
from a precomputed permutation instead of the previous load's result, which exposes
memory-level parallelism (MLP) and is expected to read \emph{faster} than the dependent
chain at the same footprint -- this is the mechanism the assignment's MLP sanity check
relies on.

\begin{lstlisting}[caption={Batched dependent and independent load timers (portable). File: \texttt{main\_code/common/benchmark.h} / \texttt{benchmark.c}, commit \texttt{e343732}.},label={lst:benchmark}]
void measure_dependency_chain_batched(struct node *start,
                                       uint64_t batch_size,
                                       size_t num_batches,
                                       double *out_latencies);

void measure_independent_loads_batched(struct node *nodes,
                                        const size_t *order,
                                        size_t order_len,
                                        uint64_t batch_size,
                                        size_t num_batches,
                                        double *out_latencies);

/* ---- benchmark.c ---- */
static volatile struct node *bench_sink;

void measure_dependency_chain_batched(struct node *start,
                                       uint64_t batch_size,
                                       size_t num_batches,
                                       double *out_latencies)
{
    struct node *p = start;

    for (size_t b = 0; b < num_batches; b++) {
        uint64_t t0 = timer_start();
        for (uint64_t i = 0; i < batch_size; i++) {
            p = p->next; /* next address depends on previous load */
        }
        uint64_t t1 = timer_stop();
        out_latencies[b] = (double)(t1 - t0) / (double)batch_size;
    }

    bench_sink = p; /* keep the final result observable */
}

void measure_independent_loads_batched(struct node *nodes,
                                        const size_t *order,
                                        size_t order_len,
                                        uint64_t batch_size,
                                        size_t num_batches,
                                        double *out_latencies)
{
    struct node *last = NULL;
    size_t idx = 0;

    for (size_t b = 0; b < num_batches; b++) {
        uint64_t t0 = timer_start();
        for (uint64_t i = 0; i < batch_size; i++) {
            /* Address comes from the caller-supplied permutation, not from
             * the value just loaded -- unlike measure_dependency_chain_batched(),
             * nothing here forces one load to wait on the previous one. */
            last = nodes[order[idx]].next;
            idx++;
            if (idx == order_len) {
                idx = 0;
            }
        }
        uint64_t t1 = timer_stop();
        out_latencies[b] = (double)(t1 - t0) / (double)batch_size;
    }

    bench_sink = last; /* keep the final result observable */
}
\end{lstlisting}

\subsubsection{Access-Pattern Construction: Randomized vs.\ Sequential Control}

Every experiment builds its dependent cycle one of two ways (Listing~\ref{lst:random}):
\texttt{make\_random\_cycle\_strided()} performs a seeded Fisher--Yates shuffle (using a
32-bit xorshift PRNG) over the node indices and links them in that shuffled order, so
consecutive accesses land at pseudo-random offsets a hardware stride prefetcher cannot
predict; \texttt{make\_sequential\_cycle\_strided()} links nodes in plain array order, a
fully stride-predictable control. Every capacity, line-size, and associativity sweep in
this project (Section~\ref{sec:phase1}) is run at \emph{both} patterns over the same swept
range: if a detected transition also appears, unchanged, on the sequential control, it is
a prefetcher or fixed-overhead artifact rather than a genuine cache-capacity/line-size/
associativity effect, and is discounted. The \emph{strided} variants generalize node
spacing beyond \texttt{sizeof(struct node)}, which is what lets the same construction be
reused unmodified for line size (spacing = candidate line size, Section~\ref{sec:linesize})
and associativity (spacing = a cache level's own confirmed capacity,
Section~\ref{sec:assoc}).

\begin{lstlisting}[caption={PRNG, shuffle, and cycle construction (portable). File: \texttt{main\_code/common/random.h} / \texttt{random.c}, commit \texttt{e343732}.},label={lst:random}]
uint32_t xorshift32(uint32_t *state) {
    uint32_t x = *state;
    x ^= x << 13;
    x ^= x >> 17;
    x ^= x << 5;
    return *state = x;
}

static struct node *node_at(void *base, size_t i, size_t stride) {
    return (struct node *)((uint8_t *)base + i * stride);
}

void make_random_cycle_strided(void *base, size_t n, size_t stride, uint32_t seed) {
    size_t *order = malloc(n * sizeof(*order));
    if (!order || n < 2) exit(1);
    for (size_t i = 0; i < n; i++) order[i] = i;

    uint32_t state = seed ? seed : 1u;
    for (size_t i = n - 1; i > 0; i--) {
        size_t j = (size_t)(xorshift32(&state) % (uint32_t)(i + 1));
        size_t tmp = order[i];
        order[i] = order[j];
        order[j] = tmp;
    }
    for (size_t i = 0; i < n; i++)
        node_at(base, order[i], stride)->next = node_at(base, order[(i + 1) % n], stride);
    free(order);
}

void make_sequential_cycle_strided(void *base, size_t n, size_t stride) {
    if (n < 2) exit(1);
    for (size_t i = 0; i < n; i++)
        node_at(base, i, stride)->next = node_at(base, (i + 1) % n, stride);
}

void make_random_cycle(struct node *nodes, size_t n, uint32_t seed) {
    make_random_cycle_strided(nodes, n, sizeof(struct node), seed);
}

void make_sequential_cycle(struct node *nodes, size_t n) {
    make_sequential_cycle_strided(nodes, n, sizeof(struct node));
}

size_t *make_shuffled_indices(size_t n, uint32_t seed) {
    size_t *order = malloc(n * sizeof(*order));
    if (!order || n < 1) exit(1);
    for (size_t i = 0; i < n; i++) order[i] = i;

    uint32_t state = seed ? seed : 1u;
    for (size_t i = n - 1; i > 0; i--) {
        size_t j = (size_t)(xorshift32(&state) % (uint32_t)(i + 1));
        size_t tmp = order[i];
        order[i] = order[j];
        order[j] = tmp;
    }
    return order;
}
\end{lstlisting}

\subsubsection{Build Configuration and the \texttt{cache\_bench} Driver}
\label{sec:build}

All six experiments are compiled with one fixed baseline (Listing~\ref{lst:makefile}):
\texttt{-O0 -g -std=c11 -Wall -Wextra -fno-omit-frame-pointer}. \texttt{-O0} is deliberate,
not just conservative -- the entire method depends on the compiler emitting one genuine
load per loop iteration in \texttt{chase()} / \texttt{measure\_dependency\_chain\_batched()}
(Listings~\ref{lst:pchase},~\ref{lst:benchmark}); any optimization pass that hoisted,
vectorized, or reordered those loads would silently break the intentional dependency the
whole technique relies on. \texttt{-g}/\texttt{-fno-omit-frame-pointer} are kept so the
emitted code can be inspected in a disassembly listing as evidence that the dependency
survives compilation -- a check flagged as still outstanding on Sunbird as of this
writing (see \texttt{data\_raw/sunbird/README.md}).

\begin{lstlisting}[style=makestyle,caption={Build configuration (\texttt{Makefile}, commit \texttt{7be6dac}).},label={lst:makefile}]
CC = gcc

CFLAGS = -O0 -g -std=c11 -Wall -Wextra -fno-omit-frame-pointer
LDLIBS = -lm

TARGET = cache_bench

COMMON_DIR = main_code/common

SOURCES = \
	$(COMMON_DIR)/main.c \
	$(COMMON_DIR)/benchmark.c \
	$(COMMON_DIR)/pointer_chase.c \
	$(COMMON_DIR)/random.c \
	$(COMMON_DIR)/capacity.c \
	$(COMMON_DIR)/line_size.c \
	$(COMMON_DIR)/associativity.c \
	$(COMMON_DIR)/latency.c \
	$(COMMON_DIR)/inclusion_policy.c

OBJECTS = $(SOURCES:.c=.o)

all: $(TARGET)

$(TARGET): $(OBJECTS)
	$(CC) $(CFLAGS) -o $@ $^ $(LDLIBS)

%.o: %.c
	$(CC) $(CFLAGS) -c $< -o $@

clean:
	rm -f $(TARGET) $(OBJECTS)

.PHONY: all clean
\end{lstlisting}

\texttt{cache\_bench} is a single binary; a \texttt{--experiment} flag selects which of the
six sweep configurations (each documented by its own header in
\texttt{main\_code/common/}, shown per-target in Section~\ref{sec:phase1}) to run, and every
common option (\texttt{--samples}, \texttt{--batch-size}, \texttt{--seed},
\texttt{--warmup-passes}, \texttt{--pattern}) threads through to whichever config struct is
active. Listing~\ref{lst:maindispatch} shows the final dispatch step; the full CLI argument
grammar and per-experiment defaults are in \texttt{main\_code/common/main.c} (commit
\texttt{7be6dac}) and the \texttt{Build \& Run} section of the repository \texttt{README.md}.

\begin{lstlisting}[caption={\texttt{cache\_bench} experiment dispatch, excerpted. File: \texttt{main\_code/common/main.c} (lines 475--503), commit \texttt{7be6dac}.},label={lst:maindispatch}]
    if (strcmp(experiment, "capacity") == 0) {
        return run_capacity_experiment(&cap_cfg);
    }
    if (strcmp(experiment, "line_size") == 0) {
        return run_line_size_experiment(&ls_cfg);
    }
    if (strcmp(experiment, "line_size_family") == 0) {
        return run_line_size_family_experiment(&fam_cfg);
    }
    if (strcmp(experiment, "associativity") == 0) {
        return run_associativity_experiment(&assoc_cfg);
    }
    if (strcmp(experiment, "hit_latency") == 0) {
        return run_hit_latency_experiment(&hl_cfg);
    }
    if (strcmp(experiment, "miss_latency") == 0) {
        return run_miss_latency_experiment(&ml_cfg);
    }
    if (strcmp(experiment, "inclusion_policy") == 0) {
        return run_inclusion_policy_experiment(&incl_cfg);
    }
\end{lstlisting}

\subsection{Machine Access, Core Pinning, and Phase Discipline}
\label{sec:phasediscipline}

Every \texttt{cache\_bench} invocation is pinned to a single physical core with
\texttt{taskset -c <core>} on the ECE lab machines, or \texttt{srun --cpu-bind=cores} on
the Hazel HPC cluster (Section~\ref{sec:hazel}). Before pinning, \texttt{Cpus\_allowed\_list}
and \texttt{lscpu -e=CPU,CORE,SOCKET,NODE} were checked to identify a full physical core
(both SMT siblings) actually available to the session, and \texttt{ps}/other-user activity
was checked to avoid a core another student was actively using. A single point-in-time
\texttt{ps} snapshot was found to be insufficient on at least one shared machine (Artemisia)
-- per-core load moved visibly across three consecutive 6-second sampling windows taken
seconds apart -- so idle-core selection there was instead based on \texttt{/proc/stat}
idle-time deltas sampled over multiple windows.

Long dense sweeps (working sets of $\gtrsim$100\,MiB at the random pattern) cost real wall
time, since random-pattern latency at that scale is dominated by genuine uncached DRAM
access; a full sweep at the default sample count can take on the order of an hour. Any run
in this range was launched inside a detached \texttt{tmux} session so it would survive the
driving terminal session disconnecting, with a second \texttt{tmux} session used purely to
poll for completion when needed -- a background-task watcher tied to an interactive session
does not itself survive that session ending, only a process actually backgrounded on the
machine (via \texttt{tmux}) does. Raw per-batch CSVs are gzip-compressed before being
committed (\texttt{data\_raw/<machine>/<experiment>/*.csv.gz}) to keep the repository size
tractable across eight machines' worth of million-sample sweeps.

One further methodological pitfall, confirmed on Sunbird, is worth stating explicitly: a
dense sweep started exactly at the left edge of the working-set range of interest can show
a spurious plateau caused by CPU frequency/P-state ramp-up (the first several points run
before the core reaches steady-state turbo clock), which decays over the first few dozen
points and is not a cache effect. Every dense sweep in this project accordingly starts
meaningfully below its region of actual interest rather than exactly at its boundary.

Finally, this project follows a strict three-phase discipline, stated in full in the
repository \texttt{README.md} and enforced by a frozen, tagged commit
(\texttt{phase1-timing-only}) before Phase II may begin:
\begin{itemize}
    \item \textbf{Phase I -- timing only} (this section and Section~\ref{sec:phase1}): no
    \texttt{perf}, no PMU, no cache-topology files, no vendor cache tables or published
    specifications used to \emph{derive} an inferred value -- every boundary, line size,
    associativity, and latency number reported in Section~\ref{sec:phase1} is inferred from
    timing measurements alone.
    \item \textbf{Phase II -- counter and literature verification} (Section~\ref{sec:pmu}):
    only entered after Phase I is tagged.
    \item \textbf{Phase III -- Hazel} (Section~\ref{sec:hazel}): a lab-only prediction is
    frozen and committed \emph{before} any cache experiment is run on Hazel.
\end{itemize}
Every experiment's exact build command, run command, compiler flags, core affinity, timer
method, sample count, seed, and raw-output filenames are recorded per machine in
\texttt{data\_raw/<machine>/README.md}, and every reduction step from raw CSV to processed
summary to plot is a checked-in script under \texttt{scripts/} (Section~\ref{sec:phase1}
lists the one used per target) -- together these are sufficient to reproduce every result
in Section~\ref{sec:phase1} from the raw data alone.

\section{Timing-Only Cache Reverse Engineering}
\label{sec:phase1}

Each subsection below follows the same chain: the access pattern and swept parameter
(built from the common code in Section~\ref{sec:build}) $\rightarrow$ the controlled sweep
$\rightarrow$ the measured distribution/plot $\rightarrow$ the inference rule applied to it
$\rightarrow$ its current Phase-II validation status. Full per-machine numeric results
(all eight machines) are tabulated in Section~\ref{sec:evolution} and in each machine's own
\texttt{data\_processed/<machine>/FINAL\_CACHE\_TABLE.md}; this section illustrates the
method itself, using Sunbird (the most complete single-machine dataset) as the running
example unless noted otherwise.

\subsection{Cache Hierarchy and Capacity}
\label{sec:capacity}

\textbf{What is varied:} working-set footprint $N$, log-spaced (8 points/octave by
default) from 1\,KiB up to a coarse ceiling (64\,MiB by default, extended manually to
1\,GiB when the topmost region has not yet plateaued). \textbf{What is held constant:}
node size, batch size, seed, warm-up passes, core pinning. \textbf{What is measured:}
median ticks/access per footprint, at both the random and sequential access patterns.
Listing~\ref{lst:capacity} is the complete implementation: for each swept size it builds a
fresh dependent cycle sized to exactly fill that footprint, runs the required warm-up
passes untimed, then times \texttt{num\_batches} batches via
\texttt{measure\_dependency\_chain\_batched()} (Listing~\ref{lst:benchmark}) and prints one
raw CSV row per (size, batch).

\begin{lstlisting}[caption={Cache-level/capacity discovery. Target property: cache-level boundaries (L1/L2/LLC/DRAM). Architecture: portable common code (compiled for both x86-64 and AArch64 via Listing~\ref{lst:timer}). File: \texttt{main\_code/common/capacity.c}, commit \texttt{8fb1467}.},label={lst:capacity}]
int run_capacity_experiment(const struct capacity_config *cfg)
{
    uint64_t min_bytes = cfg->min_bytes;
    if (min_bytes < sizeof(struct node) * 2) {
        min_bytes = sizeof(struct node) * 2;
    }

    uint64_t num_batches = (cfg->samples + cfg->batch_size - 1) / cfg->batch_size;
    if (num_batches < 1) {
        num_batches = 1;
    }
    uint64_t achieved_samples = num_batches * cfg->batch_size;

    double *batch_latencies = malloc(num_batches * sizeof(double));
    if (batch_latencies == NULL) {
        fprintf(stderr, "allocation failed for %" PRIu64 " batch samples\n", num_batches);
        return 1;
    }

    const char *pattern_name =
        (cfg->pattern == ACCESS_PATTERN_SEQUENTIAL) ? "sequential" : "random";

    printf("size_bytes,num_nodes,pattern,batch_index,avg_ticks_per_access\n");

    size_t last_num_nodes = 0;
    double step = 1.0 / (double)cfg->points_per_octave;
    double log2_min = log2((double)min_bytes);
    double log2_max = log2((double)cfg->max_bytes);

    for (double exp = log2_min; exp <= log2_max + 1e-9; exp += step) {
        uint64_t bytes = (uint64_t)(pow(2.0, exp) + 0.5);
        size_t num_nodes = (size_t)(bytes / sizeof(struct node));
        if (num_nodes < 2) num_nodes = 2;
        if (num_nodes == last_num_nodes) continue;
        last_num_nodes = num_nodes;

        struct node *nodes = malloc(num_nodes * sizeof(struct node));
        if (nodes == NULL) { free(batch_latencies); return 1; }

        if (cfg->pattern == ACCESS_PATTERN_SEQUENTIAL) {
            make_sequential_cycle(nodes, num_nodes);
        } else {
            make_random_cycle(nodes, num_nodes, cfg->seed);
        }

        /* Untimed pass(es): fault in pages and settle steady-state
         * residency for this working-set size before timing starts. */
        for (int w = 0; w < cfg->warmup_passes; w++) {
            chase(nodes, (uint64_t)num_nodes);
        }

        measure_dependency_chain_batched(nodes, cfg->batch_size, num_batches,
                                          batch_latencies);

        uint64_t size_bytes = (uint64_t)num_nodes * sizeof(struct node);
        for (uint64_t b = 0; b < num_batches; b++) {
            printf("%" PRIu64 ",%zu,%s,%" PRIu64 ",%.4f\n",
                   size_bytes, num_nodes, pattern_name, b, batch_latencies[b]);
        }
        fflush(stdout);
        free(nodes);
    }
    free(batch_latencies);
    return 0;
}
\end{lstlisting}

Raw per-batch CSVs are reduced by \texttt{scripts/summarize\_raw.py} into one
distribution-stats row per (size, pattern), then \texttt{scripts/detect\_cache\_hierarchy.py}
scans consecutive size points for a sustained tick jump exceeding both an absolute and a
relative threshold to flag a boundary, and \texttt{scripts/plot\_capacity.py} renders the
curve. \texttt{scripts/run\_capacity\_full.sh <machine> <core>} chains coarse sweep
$\rightarrow$ automatic boundary detection $\rightarrow$ per-boundary dense sweep
$\rightarrow$ reproducibility repeats $\rightarrow$ tail-extension sweep $\rightarrow$
plots into one command.

\begin{figure}[htbp]
\centering
\fbox{\parbox[c][2.2in][c]{0.85\textwidth}{\centering
PLACEHOLDER -- methodology-logic flow chart for capacity/cache-level discovery:\\[4pt]
sweep footprint $N$ (log$_2$-spaced) $\rightarrow$ build dependent cycle of $N$ bytes,
random \emph{and} sequential $\rightarrow$ batched avg.\ ticks/access $\rightarrow$
compare consecutive size points $\rightarrow$ sustained jump exceeding threshold?
$\rightarrow$ \textbf{yes}: record a cache-level boundary; \textbf{no}: continue sweep.
}}
\caption{Methodology logic for cache-level/capacity discovery. Measured quantity: median
ticks/access (y) vs.\ working-set footprint in bytes (x, log scale). Expected signature: a
sustained step in the random-pattern curve, absent (or at a different location) in the
sequential-pattern control. Inference rule: the smallest footprint where a step
exceeding both the absolute and relative thresholds is reproduced across repeats marks a
cache-level boundary.}
\label{fig:capacity-flow}
\end{figure}

\begin{figure}[htbp]
\centering
 \includegraphics[width=\textwidth]{Sunbird Capacity.png}

\caption{Capacity sweep, Sunbird (2x Intel Xeon E5-2680 v3, Haswell), core pinned via
\texttt{taskset}, 1{,}000{,}000 timed accesses/point, 8 points/octave, 1\,KiB--1\,GiB.
Units: ticks/access (TSC cycles, y-axis) vs.\ bytes (x-axis, log scale). \textbf{Boundary
values shown here are illustrative of the method only; the authoritative, hand-reconciled
per-machine L1/L2/LLC boundary values used in every downstream experiment are in
\texttt{CAPACITY\_RESULTS.md}}, not this plot in isolation. Interpretation: three
flat/ramp transitions are visible across the swept range, each confirmed via an
independent repeat run and cross-checked against the sequential-pattern control.}
\label{fig:capacity-plot}
\end{figure}

\textbf{Inference rule:} a reproducible, sustained step in median ticks/access at a
working-set size in the random-pattern curve -- confirmed not to be a frequency-ramp or
session-contention artifact (Section~\ref{sec:phasediscipline}) and not present at the same
location in the sequential-pattern control -- marks the boundary between one cache level's
occupancy and the next.

\textbf{Phase-II validation:} not yet performed on any machine; PMU-based verification
(Section~\ref{sec:pmu}) has not started for any experiment as of this writing, since Phase I
must be frozen and tagged first (Section~\ref{sec:phasediscipline}). Several machines'
automatically detected boundaries required a dense manual follow-up sweep to confirm or
retract (documented per machine in \texttt{data\_raw/<machine>/README.md}); the
consolidated, hand-verified result for every machine is \texttt{CAPACITY\_RESULTS.md}.

\subsection{Cache Line Size}
\label{sec:linesize}

\textbf{What is varied:} either the byte stride between nodes at a fixed footprint
(single-curve method), or the working-set footprint at one fixed candidate stride,
repeated once per candidate stride (family-of-curves method, PROJECT 1.pdf Figure 3 /
"Example B"). \textbf{What is held constant:} in each method, the parameter not being
swept, plus buffer base alignment (recorded, not swept, in the single-curve method; also
swept as a sub-line \texttt{offset\_bytes} cross-check in the family method).
\textbf{What is measured:} median ticks/access at each stride or footprint point, random
and sequential. Listing~\ref{lst:linesize} implements both methods in one file.

\begin{lstlisting}[caption={Cache line size, both independent methods. Target property: cache line size. Architecture: portable common code. File: \texttt{main\_code/common/line\_size.c}, commit \texttt{1689b22}.},label={lst:linesize}]
/* ---- single-curve method: fixed footprint, stride swept ---- */
int run_line_size_experiment(const struct line_size_config *cfg)
{
    uint64_t min_stride = cfg->min_stride;
    if (min_stride < sizeof(struct node)) min_stride = sizeof(struct node);

    uint64_t num_batches = (cfg->samples + cfg->batch_size - 1) / cfg->batch_size;
    if (num_batches < 1) num_batches = 1;

    double *batch_latencies = malloc(num_batches * sizeof(double));
    if (batch_latencies == NULL) return 1;

    const char *pattern_name =
        (cfg->pattern == ACCESS_PATTERN_SEQUENTIAL) ? "sequential" : "random";

    printf("stride_bytes,num_nodes,footprint_bytes,align_bytes,pattern,batch_index,avg_ticks_per_access\n");

    for (uint64_t stride = min_stride; stride <= cfg->max_stride; stride += cfg->stride_step) {
        size_t num_nodes = (size_t)(cfg->footprint_bytes / stride);
        if (num_nodes < 2) { continue; } /* stride too large for this footprint */

        uint64_t buffer_bytes = (uint64_t)(num_nodes - 1) * stride + sizeof(struct node);
        void *base = NULL;
        if (posix_memalign(&base, (size_t)cfg->align_bytes, (size_t)buffer_bytes) != 0) {
            free(batch_latencies); return 1;
        }

        if (cfg->pattern == ACCESS_PATTERN_SEQUENTIAL) {
            make_sequential_cycle_strided(base, num_nodes, (size_t)stride);
        } else {
            make_random_cycle_strided(base, num_nodes, (size_t)stride, cfg->seed);
        }

        for (int w = 0; w < cfg->warmup_passes; w++) chase(base, (uint64_t)num_nodes);
        measure_dependency_chain_batched(base, cfg->batch_size, num_batches, batch_latencies);

        for (uint64_t b = 0; b < num_batches; b++) {
            printf("%" PRIu64 ",%zu,%" PRIu64 ",%" PRIu64 ",%s,%" PRIu64 ",%.4f\n",
                   stride, num_nodes, cfg->footprint_bytes, cfg->align_bytes,
                   pattern_name, b, batch_latencies[b]);
        }
        fflush(stdout);
        free(base);
    }
    free(batch_latencies);
    return 0;
}

/* ---- family-of-curves method: fixed stride, footprint swept ----
 * Curves for stride <= true line size coincide (multiple nodes legitimately
 * share a line, so pressure depends only on total footprint); curves for
 * stride > true line size separate from that shared baseline, since each
 * node now occupies its own line. The smallest stride whose curve visibly
 * separates is the line-size estimate. See run_line_size_family_experiment()
 * in main_code/common/line_size.c for the full sweep loop (identical shape
 * to run_capacity_experiment() in Listing 8, at an arbitrary node stride
 * instead of sizeof(struct node), with an added offset_bytes cross-check). */
\end{lstlisting}

\texttt{scripts/run\_line\_size.sh <machine> <core> <boundaries\_csv>} runs the single-curve
sweep and the family-of-curves sweep (once per candidate stride, at least two
\texttt{offset\_bytes} values) per cache-level footprint taken from that machine's own
\texttt{CAPACITY\_RESULTS.md} boundary, then reports whether the two independently-designed
methods agree.

\begin{figure}[htbp]
\centering
\fbox{\parbox[c][2.0in][c]{0.85\textwidth}{\centering
PLACEHOLDER -- methodology-logic flow chart for line-size inference (both methods
side by side): \\[4pt]
\emph{Single-curve:} fix footprint just above a capacity boundary $\rightarrow$ sweep
stride linearly $\rightarrow$ ticks/access ramps then saturates $\rightarrow$ saturation
stride = line size.\\
\emph{Family-of-curves:} fix candidate stride $s$ $\rightarrow$ sweep footprint
$\rightarrow$ overlay curves for several $s$ $\rightarrow$ smallest $s$ whose curve
separates from the shared baseline = line size.
}}
\caption{Methodology logic for cache line-size inference. Expected signature: a
ramp-to-plateau elbow (single-curve) or a baseline-separation point (family-of-curves).
Inference rule: smallest stride at which the tested signature first appears.}
\label{fig:linesize-flow}
\end{figure}

\begin{figure}[htbp]
\centering
% \includegraphics[width=0.48\textwidth]{../../plots/PLACEHOLDER_linesize_single.pdf}
\includegraphics[width=0.70\textwidth]{Sunbird Linesize.png}
\caption{Line-size sweep, Sunbird L1 (32{,}768\,B footprint), both methods.
Interpretation: both independently-designed methods agree at 64\,B for this level,
confirmed with an offset cross-check.}
\label{fig:linesize-plot}
\end{figure}

\textbf{Inference rule:} single-curve -- the smallest stride at which ticks/access stops
climbing (saturates), since a larger stride beyond the true line size no longer changes how
many distinct lines are touched; family-of-curves -- the smallest candidate stride whose
footprint-vs-latency curve visibly separates from the shared low-stride baseline.

\textbf{Phase-II validation:} agreement between the two independently-designed methods
(and, for the family method, across independently repeated \texttt{offset\_bytes} values)
is treated as the internal cross-check for this target; PMU/literature comparison is
deferred to Section~\ref{sec:pmu}, not yet performed.

\subsection{Cache Associativity}
\label{sec:assoc}

\textbf{What is varied:} the number of same-set nodes chased in one dependent cycle
(\texttt{num\_ways}), from \texttt{min\_ways} to \texttt{max\_ways}. \textbf{What is held
constant:} the node-to-node stride, fixed at \texttt{cache\_bytes} -- the target cache
level's own confirmed capacity from \texttt{CAPACITY\_RESULTS.md} -- so every probed node
lands in the same cache set by construction, but carries a distinct tag.
\textbf{What is measured:} median ticks/access per \texttt{num\_ways}.

\begin{lstlisting}[caption={Cache associativity (way count per set), including the huge-page DTLB-control diagnostic. Target property: associativity. Architecture: portable common code. File: \texttt{main\_code/common/associativity.c}, commit \texttt{22104d6}.},label={lst:assoc}]
int run_associativity_experiment(const struct associativity_config *cfg)
{
    uint64_t min_ways = cfg->min_ways < 2 ? 2 : cfg->min_ways;
    uint64_t num_batches = (cfg->samples + cfg->batch_size - 1) / cfg->batch_size;
    if (num_batches < 1) num_batches = 1;

    double *batch_latencies = malloc(num_batches * sizeof(double));
    if (batch_latencies == NULL) return 1;

    const char *pattern_name =
        (cfg->pattern == ACCESS_PATTERN_SEQUENTIAL) ? "sequential" : "random";

    printf("num_ways_probed,cache_bytes,pattern,batch_index,avg_ticks_per_access\n");

    for (uint64_t num_ways = min_ways; num_ways <= cfg->max_ways; num_ways += cfg->way_step) {
        /* Every node is cfg->cache_bytes apart, so the touched address range
         * spans (num_ways - 1) strides plus one node. */
        uint64_t buffer_bytes = (num_ways - 1) * cfg->cache_bytes + sizeof(struct node);

        void *base = NULL;
        if (cfg->huge_pages) {
            /* 2 MiB-align via mmap(), madvise(MADV_HUGEPAGE), fault in every
             * 4 KiB page -- collapses the whole probe buffer onto a SINGLE
             * TLB entry regardless of num_ways, directly testing whether an
             * observed knee is a TLB/page-boundary artifact rather than a
             * real cache-set limit. Actual backing achieved is confirmed
             * from /proc/self/smaps's AnonHugePages field, not assumed from
             * a successful madvise() call. See report_huge_page_backing()
             * in main_code/common/associativity.c for the full diagnostic. */
            base = alloc_huge_page_aligned(buffer_bytes); /* see full source */
        } else {
            base = malloc((size_t)buffer_bytes);
        }
        if (base == NULL) { free(batch_latencies); return 1; }

        if (cfg->pattern == ACCESS_PATTERN_SEQUENTIAL) {
            make_sequential_cycle_strided(base, (size_t)num_ways, (size_t)cfg->cache_bytes);
        } else {
            make_random_cycle_strided(base, (size_t)num_ways, (size_t)cfg->cache_bytes, cfg->seed);
        }

        for (int w = 0; w < cfg->warmup_passes; w++) chase(base, num_ways);
        measure_dependency_chain_batched(base, cfg->batch_size, num_batches, batch_latencies);

        for (uint64_t b = 0; b < num_batches; b++) {
            printf("%" PRIu64 ",%" PRIu64 ",%s,%" PRIu64 ",%.4f\n",
                   num_ways, cfg->cache_bytes, pattern_name, b, batch_latencies[b]);
        }
        fflush(stdout);
        free_probe_buffer(base, cfg->huge_pages);
    }
    free(batch_latencies);
    return 0;
}
\end{lstlisting}

\begin{figure}[htbp]
\centering
\fbox{\parbox[c][2.0in][c]{0.85\textwidth}{\centering
PLACEHOLDER -- methodology-logic flow chart for associativity:\\[4pt]
fix stride = level's own confirmed capacity (every probed node $\rightarrow$ same set,
distinct tag) $\rightarrow$ sweep num\_ways $\rightarrow$ flat latency while
num\_ways $\le$ true associativity $\rightarrow$ sharp jump once exceeded (LRU/PLRU set
thrashes) $\rightarrow$ knee = associativity.
}}
\caption{Methodology logic for associativity inference. Expected signature: a flat
hit-latency shelf followed by a sharp, sustained jump at \texttt{num\_ways = A+1}.
Inference rule: reported associativity $A$ = the last \texttt{num\_ways} on the flat
shelf before the jump.}
\label{fig:assoc-flow}
\end{figure}

\begin{figure}[htbp]
\centering
\includegraphics[width=\textwidth]{Sunbird L1 Associativity.png}
\caption{Associativity sweep, Sunbird L1. Units: ticks/access (y) vs.\ num\_ways probed
(x). Interpretation: flat $\approx$6.0--6.2 ticks/access through \texttt{num\_ways=8},
sharp jump at 9, 0\% run-to-run spread across 3 seeded repeats $\rightarrow$ L1 = 8-way.}
\label{fig:assoc-plot}
\end{figure}

\textbf{Inference rule:} $A = $ \texttt{num\_ways} at the last flat point immediately
before the first sustained latency jump.

\textbf{Phase-II validation and an honest negative result:} L1 associativity is confirmed
cleanly this way on every machine attempted -- 8-way on three x86 machines spanning two
vendors/generations (Sunbird, Upgrade, Charnwood), 4-way on Thunderbird (AArch64,
Neoverse~N1), and a less-certain 12-way on Artemisia's larger 48\,KiB L1D. \textbf{L2 and
LLC associativity are not resolved by this method.} On every one of five machines tried,
byte-capacity candidates spanning up to a 32x range at the same nominal cache level break
at the \emph{same} small \texttt{num\_ways} (typically 9--12) regardless of the actual
byte capacity tested, and in two cases (Charnwood, Artemisia) the L2 and LLC curves are
numerically indistinguishable at every \texttt{num\_ways} despite a 32x--15x difference in
candidate capacity -- direct, on-machine evidence that a shared, capacity-independent
structure (leading hypothesis: DTLB aliasing, since every probed node lives on its own
4\,KiB page; a sliced-LLC hash function is a second, not-yet-ruled-out candidate) is what
these attempts actually measure, not the real cache level. The \texttt{--huge-pages}
control shown in Listing~\ref{lst:assoc} exists specifically to test the DTLB hypothesis by
collapsing the whole probe buffer onto one TLB entry, but has not yet been exercised on any
machine -- flagged here as the concrete next step, not attempted this cycle. A
derived-stride-scan technique (\texttt{scripts/run\_associativity\_stride\_scan.sh}) shows
more promise (self-consistent across several stride divisors on Upgrade's known-good L1)
but remains inconclusive at L2/LLC. No L2 or LLC associativity number in this report should
be read as machine-confirmed; see \texttt{CAPACITY\_RESULTS.md} and each machine's own
README for the full negative-result writeup.

\subsection{Cache Hit Latency}
\label{sec:hitlatency}

\textbf{What is varied:} load mode only (\texttt{dependent} vs.\ \texttt{independent}) at
one fixed footprint. \textbf{What is held constant:} the footprint itself, which must
already be confirmed (via Section~\ref{sec:capacity}'s boundaries) to sit entirely inside
one cache level -- this experiment infers nothing about capacity itself.
\textbf{What is measured:} median ticks/access, both load modes, both patterns.
Listing~\ref{lst:latency} implements this together with miss/next-level latency
(Section~\ref{sec:misslatency}), since both live in one file and are shown once here per
the reused-code convention.

\begin{lstlisting}[caption={Cache hit latency (this listing) and miss/next-level latency (Section~\ref{sec:misslatency} reuses this same file). Target properties: hit latency; miss/next-level latency. Architecture: portable common code. File: \texttt{main\_code/common/latency.c}, commit \texttt{e343732}.},label={lst:latency}]
int run_hit_latency_experiment(const struct hit_latency_config *cfg)
{
    size_t num_nodes = (size_t)(cfg->footprint_bytes / sizeof(struct node));
    if (num_nodes < 2) num_nodes = 2;

    uint64_t num_batches = (cfg->samples + cfg->batch_size - 1) / cfg->batch_size;
    if (num_batches < 1) num_batches = 1;

    double *batch_latencies = malloc(num_batches * sizeof(double));
    struct node *nodes = malloc(num_nodes * sizeof(struct node));
    if (batch_latencies == NULL || nodes == NULL) return 1;

    const char *pattern_name =
        (cfg->pattern == ACCESS_PATTERN_SEQUENTIAL) ? "sequential" : "random";
    const char *mode_name =
        (cfg->load_mode == LOAD_MODE_INDEPENDENT) ? "independent" : "dependent";

    if (cfg->pattern == ACCESS_PATTERN_SEQUENTIAL) {
        make_sequential_cycle(nodes, num_nodes);
    } else {
        make_random_cycle(nodes, num_nodes, cfg->seed);
    }
    /* Warm-up always via the dependent chain, regardless of the timed mode,
     * so both modes start from the same warmed-up residency. */
    for (int w = 0; w < cfg->warmup_passes; w++) chase(nodes, (uint64_t)num_nodes);

    printf("footprint_bytes,load_mode,pattern,batch_index,avg_ticks_per_access\n");

    if (cfg->load_mode == LOAD_MODE_INDEPENDENT) {
        size_t *order = (cfg->pattern == ACCESS_PATTERN_SEQUENTIAL)
            ? NULL /* identity order, built inline -- see full source */
            : make_shuffled_indices(num_nodes, cfg->seed);
        measure_independent_loads_batched(nodes, order, num_nodes, cfg->batch_size,
                                           num_batches, batch_latencies);
    } else {
        measure_dependency_chain_batched(nodes, cfg->batch_size, num_batches, batch_latencies);
    }

    for (uint64_t b = 0; b < num_batches; b++) {
        printf("%" PRIu64 ",%s,%s,%" PRIu64 ",%.4f\n",
               cfg->footprint_bytes, mode_name, pattern_name, b, batch_latencies[b]);
    }
    fflush(stdout);
    free(nodes); free(batch_latencies);
    return 0;
}

int run_miss_latency_experiment(const struct miss_latency_config *cfg)
{
    size_t target_nodes = (size_t)(cfg->target_bytes / sizeof(struct node));
    size_t evict_nodes = (size_t)(cfg->evict_bytes / sizeof(struct node));
    if (target_nodes < 2) target_nodes = 2;
    if (evict_nodes < 2) evict_nodes = 2;

    struct node *target = malloc(target_nodes * sizeof(struct node));
    struct node *evict = malloc(evict_nodes * sizeof(struct node));
    if (target == NULL || evict == NULL) return 1;

    const char *pattern_name =
        (cfg->pattern == ACCESS_PATTERN_SEQUENTIAL) ? "sequential" : "random";

    if (cfg->pattern == ACCESS_PATTERN_SEQUENTIAL) {
        make_sequential_cycle(target, target_nodes);
        make_sequential_cycle(evict, evict_nodes);
    } else {
        make_random_cycle(target, target_nodes, cfg->seed);
        make_random_cycle(evict, evict_nodes, cfg->seed + 1u);
    }
    for (int w = 0; w < cfg->warmup_passes; w++) {
        chase(target, (uint64_t)target_nodes);
        chase(evict, (uint64_t)evict_nodes);
    }

    printf("target_bytes,evict_bytes,pattern,trial_index,ticks\n");

    for (uint64_t trial = 0; trial < cfg->samples; trial++) {
        /* Untimed: re-establish target residency, then displace it by
         * walking the eviction-set cycle (evict_bytes must exceed the
         * source level's real capacity -- this is the ONLY eviction
         * mechanism in this codebase; there is no clflush anywhere). */
        chase(target, (uint64_t)target_nodes);
        chase(evict, (uint64_t)evict_nodes);

        /* Timed: exactly one dependent reload of the target line. */
        uint64_t t0 = timer_start();
        struct node *reloaded = target->next;
        uint64_t t1 = timer_stop();
        (void)reloaded;

        printf("%" PRIu64 ",%" PRIu64 ",%s,%" PRIu64 ",%" PRIu64 "\n",
               cfg->target_bytes, cfg->evict_bytes, pattern_name, trial,
               (uint64_t)(t1 - t0));
    }
    fflush(stdout);
    free(evict); free(target);
    return 0;
}
\end{lstlisting}

\begin{figure}[htbp]
\centering
\fbox{\parbox[c][1.8in][c]{0.85\textwidth}{\centering
PLACEHOLDER -- methodology-logic flow chart for hit latency:\\[4pt]
fix footprint inside one confirmed level $\rightarrow$ run BOTH dependent and
independent-load batched timers $\rightarrow$ independent reads faster (MLP exposed,
confirms the dependent timer is genuinely serializing) $\rightarrow$ report the
\emph{dependent}, random-pattern median as that level's hit latency.
}}
\caption{Methodology logic for hit-latency inference. Inference rule: dependent/random
median at a confirmed single-level footprint is the reported hit latency, valid only if
paired independent-load control reads faster at the same footprint.}
\label{fig:hitlat-flow}
\end{figure}

\begin{figure}[H]
    \centering

    % Top row
    \includegraphics[width=0.48\linewidth]{hit_latency_boxplots.pdf}
    \hfill
    \includegraphics[width=0.48\linewidth]{hit_latency_boxplots1.pdf}

    \par\medskip

    % Bottom row
    \includegraphics[width=0.48\linewidth]{hit_latency_boxplots2.pdf}
    \hfill
    \includegraphics[width=0.48\linewidth]{hit_latency_boxplots3.pdf}
    
\newpage

    \caption{Hit-latency sweep, Sunbird, footprints
    32{,}768 / 262{,}144 / 31{,}457{,}280 / 536{,}870{,}912\,B
    (L1/L2/LLC/DRAM per \texttt{CAPACITY\_RESULTS.md}).
    Units: ticks/access (TSC cycles). Interpretation: a clean,
    monotonic four-tier ladder
    (L1 $\approx 10.4$, L2 $\approx 26.8$,
    LLC $\approx 58.4$, DRAM $\approx 207$ ticks;
    dependent, random, median). The independent-load control reads
    faster than the dependent-load measurement at every level and
    pattern, confirming that MLP was exposed rather than hidden.}
    \label{fig:hitlat-plot}
\end{figure}

\subsection{Miss and Next-Level Latency}
\label{sec:misslatency}

\textbf{What is varied:} nothing within a trial -- each trial is a single timed reload;
what varies \emph{across} the sweep is which source/next-level transition is being
measured (L1$\to$L2, L2$\to$LLC, LLC$\to$DRAM), each a separate invocation.
\textbf{What is held constant:} \texttt{target\_bytes} (places the target inside the SOURCE
level) and \texttt{evict\_bytes} (an eviction-set walk sized past the source level's real
capacity but inside the next level's, per \texttt{CAPACITY\_RESULTS.md}).
\textbf{What is measured:} single-shot reload ticks per trial (not batched, since the
measurement is inherently one dependent load). The complete implementation,
\texttt{run\_miss\_latency\_experiment()}, is shown in full in Listing~\ref{lst:latency}
above (Section~\ref{sec:hitlatency}) alongside the hit-latency implementation, since both
share one source file, \texttt{main\_code/common/latency.c} (commit \texttt{e343732}); it
is not repeated here.

\begin{figure}[htbp]
\centering
\fbox{\parbox[c][1.8in][c]{0.85\textwidth}{\centering
PLACEHOLDER -- methodology-logic flow chart for miss/next-level latency:\\[4pt]
untimed re-touch of target (source-level residency restored) $\rightarrow$ untimed
eviction walk (size $>$ source capacity, $<$ next-level capacity) $\rightarrow$ ONE timed
dependent reload $\rightarrow$ record ticks (single-shot, per trial).
}}
\caption{Methodology logic for miss/next-level-latency inference. Inference rule: raw
single-shot reload latency at a transition = true reload latency + a fixed, per-machine
single-shot measurement overhead (isolated separately via a deliberately
non-colliding control eviction set, Section~\ref{sec:misslatency}).}
\label{fig:misslat-flow}
\end{figure}

\begin{figure}[H]
    \centering

    % Top-left image
    \begin{subfigure}[t]{0.48\textwidth}
        \centering
        \includegraphics[width=\linewidth]{miss_latency_boxplots.pdf}
        \caption{L1 to L2 Latentency}
        \label{fig:hitlat-a}
    \end{subfigure}
    \hfill
    % Top-right image
    \begin{subfigure}[t]{0.48\textwidth}
        \centering
        \includegraphics[width=\linewidth]{miss_latency_boxplots1.pdf}
        \caption{L2 to LLC Latency}
        \label{fig:hitlat-b}
    \end{subfigure}

    \vspace{0.5cm}

        \begin{subfigure}[t]{0.48\textwidth}
        \centering
        \includegraphics[width=\linewidth]{miss_latency_boxplots2.pdf}
        \caption{LLC to DRAM Latency}
        \label{fig:hitlat-b}
    \end{subfigure}

    \caption{Hit-latency sweep, Sunbird, footprints
    32{,}768 / 262{,}144 / 31{,}457{,}280 / 536{,}870{,}912\,B
    (L1/L2/LLC/DRAM per \texttt{CAPACITY\_RESULTS.md}).
    Units: ticks/access (TSC cycles). Interpretation: a clean,
    monotonic four-tier ladder
    (L1 $\approx 10.4$, L2 $\approx 26.8$,
    LLC $\approx 58.4$, DRAM $\approx 207$ ticks;
    dependent, random, median). The independent-load control reads
    faster than the dependent-load measurement at every level and
    pattern, confirming that MLP was exposed rather than hidden.}
    \label{fig:hitlat-plot}
\end{figure}

\textbf{Known limitation (stated directly in the source, quoted here rather than
smoothed over):} single-shot timing carries measurement overhead -- timer serialization
and post-\texttt{chase()} pipeline state cost far more when timed once than when amortized
across the $\ge$1000-access batches every other experiment uses (Section~\ref{sec:batching})
-- so a raw miss-latency median is \emph{not} directly comparable to a batched hit-latency
plateau (Section~\ref{sec:hitlatency}) without first subtracting this fixed, per-machine
overhead, which was isolated on Sunbird ($\approx$64--85 ticks) and Skylark
($\approx$65.8 ticks) via a deliberately-too-small eviction-set control but is not yet
subtracted automatically by this codebase.

\textbf{Phase-II validation:} a PMU cycle-accurate cross-check would directly resolve how
much of the raw single-shot number is fixed overhead versus real reload latency; not yet
performed.

\subsection{Inclusion and Exclusion Behavior}
\label{sec:inclusion}

\textbf{What is varied:} nothing within a trial; the sweep is across level \emph{pairings}
(e.g.\ L1 vs.\ L2, L1 vs.\ LLC, L2 vs.\ LLC). \textbf{What is held constant:}
\texttt{target\_bytes} (a small, densely-packed set confirmed inside the UPPER level) and
an untouched \texttt{control} line at the identical sub-page offset (0); the eviction walk
uses a page-multiple stride (\texttt{evict\_stride\_bytes}, a multiple of 4096) so every
eviction node shares one fixed sub-page offset (\texttt{evict\_offset\_bytes}, default
2048) different from the target/control's own offset -- this is what lets the walk
structurally avoid the target's own upper-level set while still spanning enough distinct
pages to pressure the bigger LOWER level. \textbf{What is measured:} per-trial single-shot
reload latency for both the target and the control channel.

\begin{lstlisting}[caption={Inclusion/exclusion behavior. Target property: inclusion policy between an upper (inner) and lower (outer) cache level. Architecture: portable common code. File: \texttt{main\_code/common/inclusion\_policy.c}, commit \texttt{7be6dac}.},label={lst:inclusion}]
int run_inclusion_policy_experiment(const struct inclusion_policy_config *cfg)
{
    size_t target_nodes = (size_t)(cfg->target_bytes / sizeof(struct node));
    size_t evict_nodes = (size_t)(cfg->evict_bytes / cfg->evict_stride_bytes);
    if (target_nodes < 2) target_nodes = 2;
    if (evict_nodes < 2) evict_nodes = 2;

    /* Target and control: densely packed, page-aligned (offset 0), two
     * independent buffers so control's residency can never be disturbed
     * by anything done to target's own memory. */
    struct node *target = NULL, *control = NULL;
    posix_memalign((void **)&target, 4096u, target_nodes * sizeof(struct node));
    posix_memalign((void **)&control, 4096u, target_nodes * sizeof(struct node));
    if (target == NULL || control == NULL) return 1;

    /* Eviction buffer: one node every evict_stride_bytes (a page multiple),
     * offset evict_offset_bytes past a fresh page-aligned allocation, so
     * every eviction node shares a fixed sub-page byte position different
     * from target/control's own (0) -- structurally avoids the target's own
     * set PROVIDED the upper level's index fits within one page (true for
     * a typical L1; not for an L2-or-bigger target). */
    uint64_t evict_buffer_bytes =
        (uint64_t)(evict_nodes - 1) * cfg->evict_stride_bytes + sizeof(struct node);
    uint64_t evict_alloc_bytes = evict_buffer_bytes + cfg->evict_offset_bytes;
    void *evict_alloc = NULL;
    posix_memalign(&evict_alloc, 4096u, (size_t)evict_alloc_bytes);
    if (evict_alloc == NULL) { free(control); free(target); return 1; }
    void *evict_base = (void *)((uint8_t *)evict_alloc + cfg->evict_offset_bytes);

    const char *pattern_name =
        (cfg->pattern == ACCESS_PATTERN_SEQUENTIAL) ? "sequential" : "random";

    if (cfg->pattern == ACCESS_PATTERN_SEQUENTIAL) {
        make_sequential_cycle(target, target_nodes);
        make_sequential_cycle(control, target_nodes);
        make_sequential_cycle_strided(evict_base, evict_nodes, (size_t)cfg->evict_stride_bytes);
    } else {
        make_random_cycle(target, target_nodes, cfg->seed);
        make_random_cycle(control, target_nodes, cfg->seed + 1u);
        make_random_cycle_strided(evict_base, evict_nodes, (size_t)cfg->evict_stride_bytes,
                                   cfg->seed + 2u);
    }
    for (int w = 0; w < cfg->warmup_passes; w++) {
        chase(target, (uint64_t)target_nodes);
        chase(control, (uint64_t)target_nodes);
        chase(evict_base, (uint64_t)evict_nodes);
    }

    printf("target_bytes,evict_bytes,evict_stride_bytes,evict_offset_bytes,pattern,channel,trial_index,ticks\n");

    for (uint64_t trial = 0; trial < cfg->samples; trial++) {
        /* Untimed: refresh target/control residency, displace via the
         * page-offset eviction walk (only eviction mechanism -- no clflush). */
        chase(target, (uint64_t)target_nodes);
        chase(control, (uint64_t)target_nodes);
        chase(evict_base, (uint64_t)evict_nodes);

        /* Timed: exactly one dependent reload of each. */
        uint64_t t0 = timer_start();
        struct node *r_target = target->next;
        uint64_t t1 = timer_stop();
        uint64_t t2 = timer_start();
        struct node *r_control = control->next;
        uint64_t t3 = timer_stop();
        (void)r_target; (void)r_control;

        printf("%" PRIu64 ",%" PRIu64 ",%" PRIu64 ",%" PRIu64 ",%s,target,%" PRIu64 ",%" PRIu64 "\n",
               cfg->target_bytes, cfg->evict_bytes, cfg->evict_stride_bytes,
               cfg->evict_offset_bytes, pattern_name, trial, (uint64_t)(t1 - t0));
        printf("%" PRIu64 ",%" PRIu64 ",%" PRIu64 ",%" PRIu64 ",%s,control,%" PRIu64 ",%" PRIu64 "\n",
               cfg->target_bytes, cfg->evict_bytes, cfg->evict_stride_bytes,
               cfg->evict_offset_bytes, pattern_name, trial, (uint64_t)(t3 - t2));
    }
    fflush(stdout);
    free(evict_alloc); free(control); free(target);
    return 0;
}
\end{lstlisting}

\begin{figure}[htbp]
\centering
\fbox{\parbox[c][2.1in][c]{0.85\textwidth}{\centering
PLACEHOLDER -- methodology-logic flow chart for inclusion/exclusion:\\[4pt]
touch target \& control (upper level) $\rightarrow$ page-offset eviction walk (lower
level, structurally avoids target's own set) $\rightarrow$ reload target \& control
$\rightarrow$ classify each vs.\ calibrated survived/invalidated latency classes
(Section~\ref{sec:hitlatency}/\ref{sec:misslatency}) $\rightarrow$ target invalidated
\& control survived $\Rightarrow$ inclusive; both survived $\Rightarrow$
exclusive/non-inclusive.
}}
\caption{Methodology logic for inclusion/exclusion inference. Inference rule: $\ge$80\%
of target trials classified invalidated-like $\Rightarrow$ inclusive; $\ge$80\%
survived-like $\Rightarrow$ exclusive/non-inclusive; otherwise uncertain. Control channel
must itself read survived-like on nearly every trial or the avoidance construction is not
holding for that pairing.}
\label{fig:inclusion-flow}
\end{figure}

\begin{figure}[htbp]
\centering
 \includegraphics[width=\textwidth]{inclusion_policy_boxplots.pdf}

\caption{Inclusion/exclusion, Sunbird, L1\_vs\_L2 pairing (target=L1, 32{,}768\,B;
eviction past L2, 262{,}144\,B). Units: ticks (single-shot reload). Interpretation: target
90.0\% survived-like (median 84 ticks), control 92.0\% survived-like (median 68) $\Rightarrow$
L1 is EXCLUSIVE/NON-INCLUSIVE with respect to L2, the cleanest and least caveat-exposed of
the three pairings tested on this machine. This figure illustrates the method on a single
machine; \texttt{run\_inclusion\_policy\_full.sh} has since been run for all three level
pairings on all eight lab machines (Table~\ref{tab:inclusion-crossmachine}), and the
full per-machine classifier output lives in each machine's own
\texttt{data\_processed/<machine>/FINAL\_CACHE\_TABLE.md}.}
\label{fig:inclusion-plot}
\end{figure}

\begin{table}[htbp]
\centering
\small
\caption{Inclusion/exclusion verdict by machine and level pairing, all eight lab machines
(\texttt{scripts/classify\_inclusion\_policy.py} output, one run per machine, 2026-09-13).
NI = non-inclusive/exclusive, INC = inclusive, UNC = uncertain (classifier fell short of
the 80\% threshold either way); a parenthetical lean is given for UNC verdicts where the
target/control gap still favors a direction. L2\_vs\_LLC uses an L2 target and therefore
carries the least confidence on every machine (Section~\ref{sec:inclusion}, caveat 3).}
\label{tab:inclusion-crossmachine}
\begin{tabular}{|l|l|l|l|l|}
\hline
\textbf{Machine} & \textbf{ISA/Vendor} & \textbf{L1\_vs\_L2} & \textbf{L2\_vs\_LLC} & \textbf{L1\_vs\_LLC (skip-level)} \\
\hline
Sunbird     & x86-64/Intel  & NI  & UNC (lean INC) & UNC (lean INC) \\
Artemisia   & x86-64/Intel  & UNC (footprint-confounded) & INC & INC \\
Thunderbird & AArch64/Ampere& NI  & UNC (lean INC) & NI (opposite lean from Sunbird) \\
Skylark     & x86-64/AMD    & UNC (weak lean NI) & NI  & NI \\
Upgrade     & x86-64/Intel  & NI  & NI (lean)       & UNC (lean NI) \\
Ookay       & x86-64/Intel  & NI  & UNC (lean INC) & UNC (lean INC) \\
Crux        & x86-64/Intel  & NI  & NI (lean)       & UNC (barely lean NI) \\
Charnwood   & x86-64/Intel  & NI  & UNC (lean INC) & UNC (lean INC, moderate) \\
\hline
\end{tabular}
\end{table}

\textbf{Known limitation, stated directly in the source:} the avoidance guarantee (the
eviction walk structurally cannot land on the target's own set) is proven to hold only when
the upper level's entire set-index fits within one 4\,KiB page -- true for a typical
32--64\,KiB, 8-way L1 (confirmed to fit exactly, e.g.\ Sunbird's 64 sets $\times$ 64\,B
lines = 12 index+offset bits), only \emph{approximately} true for a larger L1 with more
index bits (Thunderbird's 65{,}536\,B/4-way/256-set L1 needs 14 bits, 2 more than one page
holds), and essentially never true for an L2-or-bigger target. Pairings that use an L2
target (L2\_vs\_LLC) therefore carry the least confidence of any pairing on every machine
and are reported as leans, not firm calls. A second, independently discovered failure mode
surfaced once all eight machines were run: on Charnwood, both LLC-scale pairings' own
\emph{control} channel -- which is never touched by the eviction walk and should read
hit-speed on nearly every trial regardless of this machine's real inclusion policy --
itself showed 54.0\% and 40.0\% invalidated-like trials, direct on-machine evidence that
the large-footprint DTLB-pressure caveat (Section~\ref{sec:inclusion}'s caveat 2) can
contaminate the control channel and not just the target, something no single-machine run
had exposed. \textbf{Cross-machine summary (Table~\ref{tab:inclusion-crossmachine}):}
L1\_vs\_L2 is the one pairing that reproduces consistently -- non-inclusive/exclusive on
6 of 8 machines, with the remaining two (Artemisia, Skylark) landing UNCERTAIN rather than
disagreeing outright, one of them attributable to a known footprint-scaling confound
(Artemisia's L2 candidate is itself an unconfirmed capacity boundary). The two
LLC-involving pairings show no fleet-wide agreement at all: Artemisia reads confidently
INCLUSIVE at both, Skylark/Upgrade/Crux lean NON-INCLUSIVE at both, and
Sunbird/Thunderbird/Ookay/Charnwood land UNCERTAIN with leans that split in both
directions even between machines of the same vendor (e.g.\ Sunbird vs.\ Upgrade, both
Intel). Because the method's own structural guarantee is weakest for exactly these two
pairings, this spread is at least as consistent with the L2-target avoidance failure and
the DTLB/large-footprint confound as with genuine cross-vendor differences in LLC
inclusion policy -- this report does not claim the latter without Phase-II corroboration,
and the full per-machine reasoning (including each machine's own best-guess synthesis) is
in \texttt{data\_processed/<machine>/FINAL\_CACHE\_TABLE.md} and referenced again in
Section~\ref{sec:evolution}.

\textbf{Phase-II validation:} not yet attempted for this target on any machine.

\section{Performance Counter Verification}
\label{sec:pmu}

\subsection{Verify the Reverse Engineering with Performance Counters}

\subsubsection{Exposed Cache Events}
\texttt{perf list} confirmed real cross-vendor differences. Sunbird (Intel Haswell-EP) exposes all eight target events. Skylark (AMD Zen 2) lists only five: \texttt{L1-dcache-stores}, \texttt{LLC-loads}, and \texttt{LLC-load-misses} are absent and return \texttt{<not supported>} at runtime. Thunderbird (Arm Neoverse N1) lists six: \texttt{L1-dcache-stores} is absent, and \texttt{LLC-loads}/\texttt{LLC-load-misses} are listed but unsupported, since this PMU has no LLC-scoped event to alias them to.

\subsubsection{Counter Evidence at Boundaries}
L1 miss rate confirms the L1 boundary on all three machines: Sunbird 0.61\% to 12.59\%, Skylark 0.44\% to 9.05\%, Thunderbird 0.18\% to 2.81\%, all at the L1$\to$L2 transition. The LLC boundary is confirmed by LLC-scope miss rate on Sunbird (0.08\% to 11.25\%) and by the generic cache-miss group on Skylark (11.3\% to 46.6\%, since \texttt{LLC-loads} is unsupported there). Thunderbird has no usable last-level-miss signal at all: \texttt{LLC-loads} is unsupported, and a supplementary sweep of raw L3 events showed no capacity-dependent structure.

\subsubsection{Timing Inference vs. PMU Evidence vs. Published Value}

\begin{table}[h]
\centering
\tiny
\begin{tabular}{@{}p{2.0cm}p{2.6cm}p{2.0cm}p{2.8cm}p{1.6cm}p{3.0cm}@{}}
\toprule
Machine/Level & Timing Inference & PMU Evidence & Published/System & Source & Agreement \\
\midrule
Sunbird L1D & 32,768\,B, 8w, 64\,B, 64s, 10.35t & L1 miss 0.61\% & 32\,KiB 8w 64\,B 4cyc; sys match & Fog T10.2 & Match; latency 2.6x high \\
Sunbird L2 & 262,144\,B, 8w (guess), 64\,B, 512s, 26.83t & L1 miss 12.59\% & 256\,KiB 8w 64\,B 12cyc; sys match & Fog T10.2 & Match; guess confirmed \\
Sunbird LLC & $\sim$30\,MiB, 9w (guess), 64\,B, 58.42t & LLC miss 11.25\% & 31.46\,MB exact; 20w, 24576s; Fog 2--45\,MB 12--16w 34cyc & Fog T10.2; sysfs & Size match; ways 9 vs 20, resolved to sys \\
\midrule
Skylark L1D & 32,768\,B, 8w, 64\,B, 64s, 6.24t & L1 miss 0.44\% & 32\,KiB 8w 64\,B 4cyc; sys match & Fog T22.3 & Match; latency 1.6x high \\
Skylark L2 & 524,288\,B, 8w (guess), 64\,B, 1024s, 15.38t & L1 miss 9.05\% & 512\,KiB 8w 64\,B 12cyc; sys match & Fog T22.3 & Match; guess confirmed \\
Skylark LLC & 8.39\,MB (underestimate), 8w (guess), 128\,B, 27.26t & Cache miss 46.6\% & 16.78\,MB exact; 16w, 16384s, 64\,B; Fog 4--32\,MB 16--24w 40cyc & Fog T22.3; sysfs & Size confirms underestimate; ways 8 vs 16; line 128 vs 64 unresolved; sharing 2 cores/CCX, not socket \\
\midrule
Thunderbird L1D & 65,536\,B, 4w, 64\,B, 256s, 0.13t & L1 miss 0.20\% & 64\,KiB 4w 64\,B; sys match; no cycle figure & Ampere \S2.3.1; Arm TRM & Exact match \\
Thunderbird L2 & 1.05\,MB (taken as given), 12w (guess), 64\,B, 1365s, 0.30t & L1 miss 7.72\% & 1\,MiB 8w 64\,B 2048s; sys match & Ampere \S2.4; Arm TRM & Ways 12 vs 8, resolved by 3-way convergence \\
Thunderbird LLC & $\sim$30\,MiB, 10w (guess), 128\,B, 24576s, 0.91--1.01t & \texttt{LLC-loads} unsupported; generic 5.50\% (wrong direction) & No sysfs entry; Ampere: 32\,MiB, 16w, shared across 80 cores, mostly exclusive with L2 & Ampere \S2.5 & Size $\sim$94\%; ways 10 vs 16 unresolved; sharing wider than assumed \\
\bottomrule
\end{tabular}
\end{table}

\subsubsection{Disagreement Analysis}
Every associativity guess Phase I made above L1 disagreed with the system or literature value once checked (Sunbird's LLC: 9-way guessed vs. 20-way actual; Skylark's LLC: 8 vs. 16; Thunderbird's L2: 12 vs. 8), consistent with a shared timing confound rather than a real per-level effect. Sharing scope also surprised in both directions: Thunderbird's SLC is shared across all 80 cores rather than per socket, while Skylark's LLC is shared by only two cores per CCX rather than the whole socket. Thunderbird's inclusion result is the one place a vendor document directly contradicts Phase I: the datasheet calls the SLC mostly exclusive with L2, but the timing classifier leaned inclusive there, though that result was already flagged low-confidence. Skylark's LLC line size (128\,B measured vs. 64\,B published) remains unresolved, most likely an artifact of Zen 2's prefetcher rather than a real 128\,B line.

\Needspace{0.6\textheight}
\section{Eight Performance Counters Across Generations}

% Requires: \usepackage{booktabs}, \usepackage[T1]{fontenc}
\begin{table}[htbp]
\centering
\small
\begin{tabular}{@{} c p{4.3cm} p{9cm} @{}}
\toprule
\textbf{\#} & \textbf{Category / event name} & \textbf{Why it reflects cache theory} \\
\midrule

1 & Cache references \newline \texttt{cache-references} (same on all 3 archs) &
Denominator for a miss rate; total traffic reaching the cache hierarchy, the baseline every ratio below is measured against. \\

2 & Cache misses \newline \texttt{cache-misses} (same on all 3 archs) &
Hardware analog of the ticks/access knees from Phase I timing; should jump exactly where working-set size crosses a real capacity boundary. \\

3 & L1 data-cache accesses/loads \newline \texttt{L1-dcache-loads} (same on all 3 archs) &
Access-count baseline for the fastest, most temporally-local tier, used to normalize \#4 into a rate. \\

4 & L1 data-cache misses \newline \texttt{L1-dcache-load-misses} (same on all 3 archs) &
Direct test of locality theory: near-zero for an L1-resident chase, rising sharply once footprint exceeds L1, the counter-based confirmation of that boundary. \\

5 & L1 data-cache stores/replacements \newline Intel/AMD: \texttt{L1-dcache-stores}; ARM: \texttt{l1d\_cache\_wb} (raw; no store alias exists) &
Store traffic and eviction/writeback counts reflect how dirty lines cycle out of L1 as the chase touches new sets, which is exactly what the ``replacements'' half of the category measures. \\

6 & LLC accesses/loads \newline Intel: \texttt{LLC-loads}; AMD: inaccessible (\texttt{perf\_event\_paranoid=2} blocks the uncore PMU); ARM: \texttt{l3d\_cache} (raw) &
Traffic reaching the last level before DRAM; theoretically should stay flat once footprint exceeds L1/L2, until it itself is exceeded. \\

7 & LLC misses \newline Intel: \texttt{LLC-load-misses}; AMD: same restriction as \#6; ARM: \texttt{l3d\_cache\_refill} (raw) &
Should jump sharply once footprint exceeds LLC capacity, the counter-based confirmation of the inferred LLC boundary, mirroring \#2/\#4's logic one level up. \\

8 & dTLB misses (memory-hierarchy event) \newline \texttt{dTLB-load-misses} (same on all 3 archs) &
Tests the standing DTLB-confound hypothesis from the associativity work. If it saturates at the same footprint where the timing method's spurious num\_ways $\approx$ 9--10 wall appeared, that is hardware corroboration, not just an inferred guess. \\

\bottomrule
\end{tabular}
\caption{Eight performance events chosen for Section 8.4.}
\label{tab:8_4_events}
\end{table}

\subsection{Standardized Cross-Machine Benchmarks}
All eight events were collected on all eight ECE machines using \texttt{scripts/run\_standardized\_benchmarks.sh}, running the three standardized benchmarks (L1-resident, LLC-sized random, beyond-LLC) at each machine's own confirmed L1/LLC footprint. Two machines have structurally reduced coverage: AMD Zen 2 (Skylark) does not expose \texttt{L1-dcache-stores} or \texttt{LLC-loads/LLC-load-misses} at all in \texttt{perf list}, and ARM Neoverse N1 (Thunderbird) lacks the same two categories under \texttt{armv8\_pmuv3}. This is documented per machine rather than backfilled with a substitute.

\subsection{Normalization}
Miss rates (cache, L1, LLC) are self-normalized ratios of misses to their own reference count, so no further scaling is applied. dTLB misses are reported per 1{,}000 accesses (raw count divided by 1000, since every run fixes samples\_achieved at 1{,}000{,}000). Latency uses perf's \texttt{ns\_per\_access\_perf} rather than raw ticks, since raw tick counts are not comparable across machines with different counter frequencies (Thunderbird's 25 MHz \texttt{CNTVCT\_EL0} versus every x86 machine's multi-GHz TSC).

\subsection{Ranked Summary}

\begin{figure}[htbp]
\centering
\includegraphics[width=\textwidth]{Latency Curve.png}
\caption{Per-access latency (ns), ranked ascending within each benchmark, colored by vendor.}
\label{fig:latency_scurve}
\end{figure}

\begin{figure}
    \centering
    \includegraphics[width=\textwidth]{dtlb_miss_scurve.png}
    \caption{dTLB misses per 1{,}000 accesses, ranked ascending within each benchmark, colored by vendor.}
    \label{fig:placeholder}
\end{figure}

\begin{table}[htbp]
\centering
\small
\begin{tabular}{@{} l r r r r @{}}
\toprule
Machine (vendor, microarch) & ns/access & L1 miss \% & LLC miss \% & dTLB/1000 \\
\midrule
Crux (Intel, Coffee Lake) & 3.31 & 1.05 & 5.76 & 1.35 \\
Ookay (Intel, Kaby Lake) & 3.71 & 0.96 & 11.04 & 1.84 \\
Upgrade (Intel, Coffee Lake) & 3.95 & 0.94 & 8.82 & 1.81 \\
Charnwood (Intel, Sky Lake) & 4.27 & 1.08 & 14.00 & 2.03 \\
Skylark (AMD, Zen 2) & 5.33 & 0.44 & n/a & 0.69 \\
Thunderbird (Ampere/ARM, N1) & 5.84 & 0.19 & n/a & 1.56 \\
Artemisia (Intel, Sapphire Rapids) & 6.27 & 0.61 & 8.69 & 1.69 \\
Sunbird (Intel, Haswell) & 10.19 & 0.63 & 2.71 & 0.51 \\
\bottomrule
\end{tabular}
\caption{L1-resident benchmark, ranked ascending by latency.}
\label{tab:l1_resident}
\end{table}

\begin{table}[htbp]
\centering
\small
\begin{tabular}{@{} l r r r r @{}}
\toprule
Machine (vendor, microarch) & ns/access & L1 miss \% & LLC miss \% & dTLB/1000 \\
\midrule
Skylark (AMD, Zen 2) & 74.04 & 6.13 & n/a & 2.91 \\
Crux (Intel, Coffee Lake) & 88.03 & 8.80 & 3.59 & 1080.28 \\
Ookay (Intel, Kaby Lake) & 136.89 & 8.79 & 16.33 & 1091.71 \\
Charnwood (Intel, Sky Lake) & 174.84 & 8.83 & 15.92 & 1096.98 \\
Upgrade (Intel, Coffee Lake) & 203.70 & 9.52 & 18.75 & 3091.93 \\
Sunbird (Intel, Haswell) & 542.14 & 7.54 & 7.25 & 0.92 \\
Artemisia (Intel, Sapphire Rapids) & 1183.75 & 8.76 & 30.68 & 11084.29 \\
Thunderbird (Ampere/ARM, N1) & 1296.89 & 5.46 & n/a & 1638.88 \\
\bottomrule
\end{tabular}
\caption{LLC-sized random benchmark, ranked ascending by latency.}
\label{tab:llc_random}
\end{table}

\begin{table}[htbp]
\centering
\small
\begin{tabular}{@{} l r r r r @{}}
\toprule
Machine (vendor, microarch) & ns/access & L1 miss \% & LLC miss \% & dTLB/1000 \\
\midrule
Crux (Intel, Coffee Lake) & 19530.32 & 13.42 & 46.75 & 254971.01 \\
Upgrade (Intel, Coffee Lake) & 19788.38 & 13.36 & 47.05 & 256121.64 \\
Ookay (Intel, Kaby Lake) & 19847.98 & 13.36 & 47.79 & 252964.81 \\
Thunderbird (Ampere/ARM, N1) & 25154.25 & 5.34 & n/a & 255722.15 \\
Charnwood (Intel, Sky Lake) & 28545.55 & 13.42 & 47.53 & 253582.41 \\
Skylark (AMD, Zen 2) & 29989.45 & 5.52 & n/a & 10.90 \\
Artemisia (Intel, Sapphire Rapids) & 35782.53 & 11.21 & 76.72 & 251331.33 \\
Sunbird (Intel, Haswell) & 38588.18 & 12.94 & 49.48 & 210666.08 \\
\bottomrule
\end{tabular}
\caption{Beyond-LLC benchmark, ranked ascending by latency.}
\label{tab:beyond_llc}
\end{table}

\subsection{Cross-Vendor and Cross-Generation Comparison}
LLC-scope counters and L1 store counts are not comparable across all three vendors: both are absent on AMD and ARM, leaving only the six Intel machines with that data. The generic \texttt{cache-references}/\texttt{cache-misses} pair is present everywhere but is not semantically equivalent across PMUs either, since it has already been shown to track L1-scope traffic almost exactly on Thunderbird's ARM PMU rather than genuine last-level activity, so raw comparisons of that pair across vendors would be misleading. dTLB-load-misses is the most portable of the eight events, with consistent semantics on all three vendors.

Skylark (AMD) shows a striking and reproducible dTLB anomaly: at LLC-sized and beyond-LLC footprints its dTLB miss rate is two to four orders of magnitude lower than every other machine at the same tier (2.91 and 10.90 per 1{,}000 accesses versus 1000+ and 250{,}000+ elsewhere), suggesting a real architectural difference in TLB reach or page management rather than measurement noise, since it is consistent across both larger benchmarks. Within Intel alone, the oldest machine (Sunbird, Haswell, 2014) shows the highest L1-resident and beyond-LLC latency of the Intel group, while the newest (Artemisia, Sapphire Rapids, 2023) is a partial outlier with the highest LLC-random latency and dTLB miss count among Intel machines, most likely reflecting its much larger multi-socket server die rather than a simple generational speedup, so raw latency does not track chip age monotonically once die size and topology differ this much.

\section{Software-Only Cache Metric}

\section{Chronological Cache Evolution Analysis}
\label{sec:evolution}

\section{Cache Evolution Prediction}

\subsection{Your Team's Cache Laws and Future Scaling Walls}
\label{sec:cachelaws}

Two named laws are formulated below, using both team members' surnames (Chen, Ngo) per the
assignment's naming convention. Both are fit exclusively from Phase I, timing-only measurements on
the team's 8 ECE lab machines (\texttt{CHRONOLOGICAL\_MASTER\_TABLE.md}; backing data
\texttt{data\_processed/master/chronological\_master\_table.csv}) -- no Hazel data was used to fit
either law. They concern two structurally different cache-related quantities, per the assignment's
requirement: per-core L2 capacity (a growth law) and L1D capacity/associativity/line size (an
approximate-constancy law). Both laws' parameters and their Hazel predictions were frozen together in
\texttt{PREDICTION\_FREEZE.md} at commit \texttt{b7d98c0} (git tag \texttt{predictions-frozen},
2026-09-14 18:23 UTC) \textbf{before} any Hazel cache experiment ran, satisfying this project's Phase
Discipline (\S\ref{sec:phasediscipline}).

\subsubsection*{Law 1: Chen--Ngo Cache Capacity Doubling Law}

\textbf{Statement.} Per-core L2 cache capacity on the team's lab machines grows with processor
introduction year via a step-doubling trend, doubling on average every 3 years across the 2014--2023
span sampled.
\noindent
\textbf{Quantitative rule.} $C(t) = 256\,\text{KiB} \cdot 2^{(t - 2014)/\tau}$, with $\tau = 3$ years,
fit from the dataset's two chronological endpoints (Sunbird, 2014, 256\,KiB $\to$ Artemisia, 2023,
2\,MiB -- an exact $8\times$ increase over 9 years). The underlying data is a step function, not a
smooth curve: flat at 256\,KiB for five years (2014--2019: Sunbird, Charnwood, Ookay, Upgrade, Crux),
then three back-to-back doublings in four years (512\,KiB -- Skylark, 2019; 1\,MiB -- Thunderbird,
2020; 2\,MiB -- Artemisia, 2023). The fitted $\tau$ is an end-to-end summary between two endpoints, not
a per-year validated slope (see \S\ref{sec:evolution}, items 2, 5, 6, for the full derivation).
\noindent
\textbf{Scope/domain.} 8 machines, 2014--2023, pooled across vendor/ISA (6 Intel x86-64, 1 AMD
x86-64 [Skylark], 1 Arm AArch64 [Thunderbird]) and pooled across desktop/server class. Not
vendor-specific: the two fitted endpoints are both Intel, but the two intermediate step points come
from AMD and Arm parts, so the doubling time is not read off a single vendor's roadmap alone. Does not
imply the trend holds beyond 2023, outside this 8-machine sample, or for a machine class not
represented in this fleet.

\begin{figure}[H]
\centering
\fbox{
\begin{minipage}[c][7cm][c]{0.92\textwidth}
\centering
\textbf{\large PLOT PLACEHOLDER --- TODO}

\vspace{10pt}

\itshape
X-axis: processor introduction year (2014--2028). Y-axis: per-core L2 capacity, log$_2$(KiB).
Solid points + solid line: the 8 lab machines, oldest (Sunbird, 2014) to newest (Artemisia, 2023),
per \texttt{CHRONOLOGICAL\_MASTER\_TABLE.md}'s L2 table. Dashed line: the fitted model
$C(t) = 256\,\text{KiB}\cdot 2^{(t-2014)/3}$ extrapolated forward to $\sim$2028 and evaluated at
Hazel's \texttt{haswell} target generation ($t \approx 2014$).
\end{minipage}
}

\caption{Evidence plot for the Chen--Ngo Cache Capacity Doubling Law -- at least 3 lab-machine
observations as solid points/line, model extrapolation as a dashed line, Moore-style format.}
\label{fig:law1-placeholder}
\end{figure}
\noindent
\textbf{Frozen prediction.} Evaluated at Hazel's target \texttt{haswell} generation ($t \approx 2014$,
the model's own lower fit endpoint): \textbf{256\,KiB} per core -- a same-era identity prediction, not
a forward extrapolation, so it carries no extra uncertainty band beyond the step-function caveat
above. Full frozen record: \texttt{PREDICTION\_FREEZE.md}, ``Team Cache Laws -- Frozen Predictions''
table.

\subsubsection*{Law 2: Chen--Ngo Cache Invariance Law}

\textbf{Statement.} L1D cache capacity, associativity, and line size on the team's lab machines are
approximately invariant across a decade of processor generations and three vendors/ISAs, rather than
scaling with year the way L2 capacity does.
\noindent
\textbf{Quantitative rule.} Line size: 64\,B on 8 of 8 machines (2014--2023), zero exceptions. L1D
capacity: 32,768\,B (32\,KiB) on 6 of 8 machines. L1D associativity: 8-way on 6 of 8 machines. The
same two machines account for both deviations -- Thunderbird (Arm, 2020: 64\,KiB, 4-way) and Artemisia
(Intel Sapphire Rapids, 2023: 48\,KiB, 12-way) -- and they deviate in different directions/magnitudes
from each other despite Artemisia being the newer machine, i.e.\ the deviations themselves show no
chronological trend.
\noindent
\textbf{Scope/domain.} Same 8-machine, 2014--2023, cross-vendor/ISA pool as Law 1. The invariance
claim is majority-case (6 of 8), not universal, and is stated as approximate; it is known not to hold
for the team's one Arm machine or its newest/largest-L1D Intel Xeon, and should not be assumed to hold
for machine classes outside this fleet.

\begin{figure}[H]
\centering
\fbox{
\begin{minipage}[c][7cm][c]{0.92\textwidth}
\centering
\textbf{\large PLOT PLACEHOLDER --- TODO}

\vspace{10pt}

\itshape
X-axis: processor introduction year (2014--2023). Y-axis: L1D capacity in KiB (left) and line size
in bytes (right, or a second panel). Solid points: the 8 lab machines --- 6 clustered flat at
32\,KiB/64\,B, with Thunderbird (2020, 64\,KiB) and Artemisia (2023, 48\,KiB) marked as labeled
outliers. Dashed line: the constant-value prediction (32\,KiB / 8-way / 64\,B) held forward to
Hazel's \texttt{haswell} target generation.
\end{minipage}
}
\caption{Evidence plot for the Chen--Ngo Cache Invariance Law -- at least 3 lab-machine observations
as solid points, the constant-value prediction as a dashed line.}
\label{fig:law2-placeholder}
\end{figure}

\noindent
\textbf{Frozen prediction.} At Hazel's \texttt{haswell} target: L1D capacity \textbf{32,768\,B},
associativity \textbf{8-way}, line size \textbf{64\,B} -- chosen because \texttt{haswell} is an x86
Xeon target, the exact class where this law's 6/8 majority held with no exceptions (both deviations
are non-Haswell parts). Full frozen record: \texttt{PREDICTION\_FREEZE.md}.

\subsubsection*{Future Scaling Walls}

\textbf{Wall for Law 1 (capacity doubling).} The 3-year doubling time is fit from a fleet that was flat
for 5 years and then doubled three times in 4 years -- i.e.\ the model is extrapolating a recent burst,
not a long-observed steady rate, and two independent physical constraints could break it going
forward. First, a \emph{latency wall}: \S\ref{sec:evolution} item 3 already shows LLC hit latency
broadly growing with LLC capacity across this fleet; the same wire-delay/capacity tradeoff applies to
L2, so continuing to double L2 capacity every 3 years would eventually push L2 hit latency toward
today's LLC latency range, eroding the two-level (L2/LLC) split's purpose unless vendors insert an
additional mid-level cache -- a redesign, not a continuation of the same curve. Second, a
\emph{fabrication wall}: this project's own data cannot observe process-node effects directly (Phase I
is timing-only), but the fitted doublings land almost entirely in the fleet's more recent, smaller-node
machines (Skylark and Thunderbird at 7\,nm, Artemisia at 10\,nm/Intel~7), while SRAM cell density is
publicly known to scale much more slowly than logic density below roughly 7\,nm -- a hard limit this
project's data is consistent with but cannot independently confirm, flagged here as an external risk
to the extrapolation rather than a lab finding. Third, a \emph{core-count wall}: per-core L2 capacity
competes with core count for the same fixed die-area/power budget (already noted as a risk in
\S\ref{sec:evolution} item 7); if future generations grow core counts faster than die area, per-core L2
capacity could flatten or reverse even while total on-die SRAM keeps growing.

\medskip

\noindent\textbf{Wall for Law 2 (invariance).} Unlike Law 1, this law describes a quantity that appears to have
\emph{already} hit its scaling wall rather than one that will hit it in the future: L1D capacity and
associativity are bounded by a strict load-to-use latency budget (commonly $\leq$4 cycles in modern
designs), and growing L1D physically increases wire delay and tag-comparator fan-in in a way that would
violate that budget -- consistent with why L1D has not grown across 9 years and 3 vendors on 6 of 8
machines even as transistor budgets elsewhere on the same dies have grown by orders of magnitude. The
two deviations show the wall is soft, not absolute: a different microarchitectural latency budget
(Arm's Neoverse~N1 core design; Intel's newest server core) can push through it by roughly
50--100\,\%, but not by the multiplicative factors Law 1's L2 trend shows over the same span. The
practical implication for future generations, including any that show a materially larger L1D than
32\,KiB, is that this project's own timing-only method predicts it should also show a
correspondingly \emph{higher} L1D hit-latency in ticks (not a free capacity increase) -- a testable,
falsifiable consequence of this law that a future Hazel or successor-fleet measurement could check
directly against the frozen 32,768\,B / 8-way / 64\,B prediction above.

The held-out evaluation for both laws (measured Hazel value, error, and supported/weakened/falsified
verdict) is recorded in \texttt{PREDICTION\_FREEZE.md}'s ``Held-Out Evaluation'' table once the Hazel
cache-experiment suite has run; per that file's own rule, Hazel points are never used to refit either
law above.

Item 7's forward-looking capacity prediction (the general Section~\ref{sec:evolution} extrapolation,
distinct from the two named laws above) is derived and stated in full in Section~\ref{sec:evolution}
above, alongside the rest of the Section 9 interpretation items, since it follows directly from the
same doubling-time model (item 6). The held-out evaluation half of item 7 is likewise pending Hazel
data and will be added to that item once Hazel measurements are collected.

\section{Hazel HPC Validation}
\label{sec:hazel}

\section{Discussion}

\section{Conclusion}

\newpage

% ============================================================
% MANDATORY TEAM CONTRIBUTION APPENDIX
% ============================================================

% Force the next section to be Section 16.
% Do NOT use \appendix because that would make the section
% number change from 16 to A.
\setcounter{section}{12}
\renewcommand{\thesection}{\arabic{section}}

% Make subsections use letters A, B, C, ..., I.
\setcounter{subsection}{0}
\renewcommand{\thesubsection}{\Alph{subsection}}

\section{Mandatory Team Contribution Appendix}

\subsection{Team Members and GitHub Accounts}

\begin{table}[htbp]
\centering
\caption{Team-member identification and GitHub accounts for the mandatory contribution appendix.}
\label{tab:team-members}

\begin{tabular}{|p{4cm}|p{4cm}|p{5cm}|}
\hline
\textbf{Member Name} &
\textbf{Unity ID} &
\textbf{GitHub Username} \\
\hline

Member 1 &
nsngo &
NicholasNgo4 \\
\hline

Member 2 &
krchen &
krchen1 \\
\hline

\end{tabular}

\end{table}


\subsection{Result / Data-Point Attribution}

For every result shown in the paper---including each final table entry, plotted machine/generation point, prediction point, PMU comparison, and major derived quantity---identify who produced the underlying measurement and who checked/processed it.

Mechanically repeated points from one identical sweep may be grouped as one row. For example, ``L1 capacity sweep on Sunbird, all working-set sizes.'' However, no reported result should be unattributed.

\begin{longtable}{|p{1.2cm}|p{4cm}|p{2.5cm}|p{2.5cm}|p{3cm}|}
\caption{Result/data-point attribution for every reported result.}
\label{tab:result-attribution} \\

\hline
\textbf{Result / Data ID} &
\textbf{Result, figure, table, machine, or data series} &
\textbf{Member 1 contribution} &
\textbf{Member 2 contribution} &
\textbf{Raw-data / script / commit reference} \\
\hline
\endfirsthead

\hline
\textbf{Result / Data ID} &
\textbf{Result, figure, table, machine, or data series} &
\textbf{Member 1 contribution} &
\textbf{Member 2 contribution} &
\textbf{Raw-data / script / commit reference} \\
\hline
\endhead

R1 & & & & \\
\hline
R2 & & & & \\
\hline
R3 & & & & \\
\hline
R4 & & & & \\
\hline
R5 & & & & \\
\hline
R6 & & & & \\
\hline
R7 & & & & \\
\hline
R8 & & & & \\
\hline
R9 & & & & \\
\hline
R10 & & & & \\
\hline
R11 & & & & \\
\hline
R12 & & & & \\
\hline

\end{longtable}

Add rows as needed. Your report may require substantially more than 12 rows.


\subsection{Team Meeting Log}

Record project meetings or substantial synchronous work sessions, if any. Include who attended and a concise summary of decisions, work completed, unresolved issues, and next steps.

\begin{longtable}{|p{2.5cm}|p{3cm}|p{2cm}|p{2cm}|p{5cm}|}
\caption{Team meeting and synchronous-work log.}
\label{tab:meeting-log} \\

\hline
\textbf{Date / Time} &
\textbf{Purpose} &
\textbf{Member 1 attended?} &
\textbf{Member 2 attended?} &
\textbf{Summary / decisions / next steps} \\
\hline
\endfirsthead

\hline
\textbf{Date / Time} &
\textbf{Purpose} &
\textbf{Member 1 attended?} &
\textbf{Member 2 attended?} &
\textbf{Summary / decisions / next steps} \\
\hline
\endhead

29Aug2026, 3pm &
Meet Team Members, start on project &
Yes &
Yes &
Github + overleaf report created, tested access to machines, filled out Table 1, begin building micro benchmark\\
\hline

& & & & \\
\hline

& & & & \\
\hline

& & & & \\
\hline

\end{longtable}

Add rows as needed. If the team had no formal meetings, state that explicitly and describe how coordination occurred.


\subsection{Code Contribution}

List the code components, scripts, benchmark modules, build files, and infrastructure each member created, modified, reviewed, or debugged. Include Git commit hashes or pull-request references where useful.

\begin{longtable}{|p{4cm}|p{3.5cm}|p{3.5cm}|p{3cm}|}
\caption{Code, script, build, and infrastructure contributions.}
\label{tab:code-contributions} \\

\hline
\textbf{Code component / file / module} &
\textbf{Member 1 contribution} &
\textbf{Member 2 contribution} &
\textbf{Git reference} \\
\hline
\endfirsthead

\hline
\textbf{Code component / file / module} &
\textbf{Member 1 contribution} &
\textbf{Member 2 contribution} &
\textbf{Git reference} \\
\hline
\endhead

& & & \\
\hline

& & & \\
\hline

& & & \\
\hline

& & & \\
\hline

& & & \\
\hline

\end{longtable}


\subsubsection*{Required GitHub Contribution-Plot Screenshot}

Immediately after the code-contribution table, insert a screenshot of the GitHub repository's contributor/contribution plot.

\begin{figure}[htbp]
\centering

% Uncomment after adding the screenshot:
% \includegraphics[width=\textwidth]{github_contributions.png}

\fbox{
\parbox[c][3in][c]{0.85\textwidth}{
\centering
Insert GitHub contributor/contribution plot screenshot here.

Show both student accounts; caption with full name and GitHub username mapping.
}
}

\caption{GitHub contributor/contribution evidence.}
\label{fig:github-contributions}

\end{figure}


\subsection{Experimental-Running Contribution}

Document who actually ran each major experiment. Include the machine/generation and enough information to connect the work to the raw-data files.

\begin{longtable}{|p{4.5cm}|p{3cm}|p{4cm}|p{4cm}|}
\caption{Experimental-running contribution by machine/generation.}
\label{tab:experimental-contributions} \\

\hline
\textbf{Experiment / machine / generation} &
\textbf{Run date(s)} &
\textbf{Member 1 contribution} &
\textbf{Member 2 contribution} \\
\hline
\endfirsthead

\hline
\textbf{Experiment / machine / generation} &
\textbf{Run date(s)} &
\textbf{Member 1 contribution} &
\textbf{Member 2 contribution} \\
\hline
\endhead

& & & \\
\hline

& & & \\
\hline

& & & \\
\hline

& & & \\
\hline

& & & \\
\hline

\end{longtable}


\subsection{Plotting and Writing Contribution}

List each major figure/table and each major paper section. State who generated/checked the plots, who drafted the text, and who revised it.

\begin{longtable}{|p{5cm}|p{4.5cm}|p{4.5cm}|}
\caption{Plotting, table, and report-writing contributions.}
\label{tab:plotting-writing-contributions} \\

\hline
\textbf{Figure / table / report section} &
\textbf{Member 1 contribution} &
\textbf{Member 2 contribution} \\
\hline
\endfirsthead

\hline
\textbf{Figure / table / report section} &
\textbf{Member 1 contribution} &
\textbf{Member 2 contribution} \\
\hline
\endhead

& & \\
\hline

& & \\
\hline

& & \\
\hline

& & \\
\hline

& & \\
\hline

\end{longtable}


\subsection{Analysis Contribution}

Document the reasoning work separately from running experiments and making plots.

\begin{longtable}{|p{5cm}|p{4.5cm}|p{4.5cm}|}
\caption{Analysis and interpretation contributions.}
\label{tab:analysis-contributions} \\

\hline
\textbf{Analysis question / conclusion} &
\textbf{Member 1 contribution} &
\textbf{Member 2 contribution} \\
\hline
\endfirsthead

\hline
\textbf{Analysis question / conclusion} &
\textbf{Member 1 contribution} &
\textbf{Member 2 contribution} \\
\hline
\endhead

& & \\
\hline

& & \\
\hline

& & \\
\hline

& & \\
\hline

& & \\
\hline

\end{longtable}


\subsection{AI / LLM Assistance Disclosure}

Complete this table if any AI/LLM tool was used. Add rows as needed. If no AI was used, write ``None.''

\begin{longtable}{|p{3cm}|p{3.5cm}|p{4cm}|p{4.5cm}|}
\caption{AI/LLM assistance disclosure.}
\label{tab:ai-disclosure} \\

\hline
\textbf{Tool} &
\textbf{Task / artifact} &
\textbf{Type of help received} &
\textbf{What the team verified and can explain} \\
\hline
\endfirsthead

\hline
\textbf{Tool} &
\textbf{Task / artifact} &
\textbf{Type of help received} &
\textbf{What the team verified and can explain} \\
\hline
\endhead

Claude &
GitHub Structure and Necessary Folders &
Content Creation - Created Github Infrastructure containing required folders/files &
The content creation is purely administrative. The populated folders and files did not contain any technical content and were created based off what the project specifications required. The team verified the repository by examining the contents inside it.\\
\hline

ChatGPT &
Overleaf LaTeX Report
& Content Creation - Populated report section headers and fixed syntax issues
& The team verified by examining the LaTeX file, ensuring the file compiles, and content matches the project requirements.\\
\hline

Claude &
Report Sections 3--4 (Experimental Methodology; Timing-Only Cache Reverse Engineering) &
Content Creation -- drafted methodology prose, inserted verified source-code listings copied verbatim from \texttt{main\_code/common/} with file/commit citations, and added methodology-logic/plot figure placeholders &
The team verified every code listing against the actual repository source at the cited commit, confirmed the described method matches each header's own docstring and each machine's \texttt{data\_raw/<machine>/README.md}, and can explain the pointer-chase/timing/batching/access-pattern design and every stated caveat (e.g.\ the unresolved L2/LLC associativity confound, the single-shot miss-latency overhead) independently of the AI-drafted text.\\
\hline

\end{longtable}

\noindent
\textbf{AI accountability statement:} Both students are responsible for understanding every submitted artifact, including AI-assisted code and analysis.


\subsection{Final Contribution Confirmation}

Both members must review the completed appendix before submission and complete Table 14.


\begin{table}[H]
\centering
\caption{Final contribution-appendix confirmation by both team members.}
\label{tab:final-confirmation}

\begin{tabular}{|p{6cm}|p{8cm}|}
\hline
\textbf{Member Name} &
\textbf{Confirmation / Date} \\
\hline

Member 1 & \\
\hline

Member 2 & \\
\hline

\end{tabular}

\end{table}

\end{document}
